\input{preamble}

% OK, start here.
%
\begin{document}

\title{Espaces algébriques sur les corps}


\maketitle

\phantomsection
\label{section-phantom}

\tableofcontents

\section{Introduction}
\label{section-introduction}

\noindent
Ce chapitre est l'analogue du chapitre sur les variétés dans le cadre
des espaces algébriques. Une référence pour les espaces algébriques est
\cite{Kn}.


\section{Conventions}
\label{section-conventions}

\noindent
Nous faisons l'hypothèse permanente que tous les schémas appartiennent à
un gros site fppf $\Sch_{fppf}$. De plus, tout anneau $A$ considéré
est tel que $\Spec(A)$ est (isomorphe à) un
objet de ce grand site.

\medskip\noindent
Soit $S$ un schéma et soit $X$ un espace algébrique sur $S$.
Dans ce chapitre et le suivant, nous écrirons $X \times_S X$
pour le produit de $X$ par lui-même (dans la catégorie des espaces algébriques
sur $S$), au lieu de $X \times X$.



\section{Morphismes génériquement finis}
\label{section-generically-finite}

\noindent
Cette section prolonge la discussion de
Espaces décents, section \ref{decent-spaces-section-generically-finite},
ainsi que son analogue, pour les morphismes d'espaces algébriques, dans
Variétés, section \ref{varieties-section-generically-finite}.

\begin{lemma}
\label{lemma-quasi-finite-in-codim-1}
Soit $S$ un schéma. Soit $f : X \to Y$ un morphisme d'espaces algébriques
sur $S$. Supposons que $f$ soit localement de type fini et que $Y$ soit
localement noethérien. Soit $y \in |Y|$ un point de codimension $\leq 1$ de $Y$.
Soit $X^0 \subset |X|$ l'ensemble des points de codimension $0$ de $X$.
Supposons de plus que l'une des conditions suivantes soit satisfaite
\begin{enumerate}
\item pour tout $x \in X^0$, le degré de transcendance de $x/f(x)$ est $0$,
\item pour tout $x \in X^0$ tel que $f(x) \leadsto y$,
le degré de transcendance de $x/f(x)$ est $0$,
\item $f$ est quasi-fini en tout $x \in X^0$,
\item $f$ est quasi-fini en tout point d'une partie dense de $|X|$,
\item en ajouter d'autres ici.
\end{enumerate}
Alors $f$ est quasi-fini en tout point de $X$ situé au-dessus de $y$.
\end{lemma}

\begin{proof}
Nous voulons ramener la preuve au cas des schémas. Pour ce faire, nous
choisissons un diagramme commutatif
$$
\xymatrix{
U \ar[r] \ar[d]_g & X \ar[d]^f \\
V \ar[r] & Y
}
$$
où $U$ et $V$ sont des schémas et où les flèches horizontales sont \'etales
et surjectives. Choisissons $v \in V$ d'image $y$. Observons que
$V$ est localement noethérien et que $\dim(\mathcal{O}_{V, v}) \leq 1$
(voir Propriétés des espaces, Définitions
\ref{spaces-properties-definition-dimension-local-ring} et
Remarque \ref{spaces-properties-remark-list-properties-local-etale-topology}).
La fibre $U_v$ de $U \to V$ au-dessus de $v$ s'envoie surjectivement sur
$f^{-1}(\{y\}) \subset |X|$. L'image réciproque de $X^0$ dans $U$
est exactement l'ensemble des
points génériques des composantes irréductibles de $U$ (Propriétés des espaces, Lemme
\ref{spaces-properties-lemma-codimension-0-points}).
Si $\eta \in U$ est un tel point, d'image $x \in X^0$, alors
le degré de transcendance de $x / f(x)$ est le degré de transcendance de
$\kappa(\eta)$ sur $\kappa(g(\eta))$
(Morphismes d'espaces, Définition
\ref{spaces-morphisms-definition-dimension-fibre}).
Observons que $U \to V$ est quasi-fini en $u \in U$ si et seulement si
$f$ est quasi-fini en l'image de $u$ dans $X$.

\medskip\noindent
Cas (1). Ici le cas (1) de
Variétés, Lemme \ref{varieties-lemma-quasi-finite-in-codim-1} s'applique
et nous concluons que $U \to V$ est quasi-fini en tout point de $U_v$.
Donc $f$ est quasi-fini en tout point situé au-dessus de $y$.

\medskip\noindent
Cas (2). Soit $u \in U$ un point générique d'une composante irréductible
dont l'image dans $V$ se spécialise en $v$. Alors l'image $x \in X^0$ de
$u$ vérifie $f(x) \leadsto y$. Par conséquent, nous voyons que
le cas (2) de
Variétés, Lemme \ref{varieties-lemma-quasi-finite-in-codim-1} s'applique
et nous concluons comme précédemment.

\medskip\noindent
Le cas (3) découle du cas (3) de
Variétés, Lemme \ref{varieties-lemma-quasi-finite-in-codim-1}.

\medskip\noindent
Dans le cas (4), puisque l'application $|U| \to |X|$ est ouverte, nous voyons que
l'ensemble des points où $U \to V$ est quasi-fini est également dense.
Ainsi le cas (4) de
Variétés, Lemme \ref{varieties-lemma-quasi-finite-in-codim-1} s'applique.
\end{proof}

\begin{lemma}
\label{lemma-finite-in-codim-1}
Soit $S$ un schéma. Soit $f : X \to Y$ un morphisme d'espaces algébriques
sur $S$. Supposons que $f$ soit propre et que $Y$ soit localement
noethérien. Soit $y \in Y$ un point de codimension $\leq 1$ de $Y$.
Soit $X^0 \subset |X|$ l'ensemble des points de codimension $0$ de $X$.
Supposons de plus que l'une des
conditions suivantes soit satisfaite
\begin{enumerate}
\item pour tout $x \in X^0$, le degré de transcendance de $x/f(x)$ est $0$,
\item pour tout $x \in X^0$ tel que $f(x) \leadsto y$, le degré de transcendance
de $x/f(x)$ est $0$,
\item $f$ est quasi-fini en tout $x \in X^0$,
\item $f$ est quasi-fini en tout point d'une partie dense de $|X|$,
\item en ajouter d'autres ici.
\end{enumerate}
Alors il existe un sous-espace ouvert $Y' \subset Y$ contenant $y$ tel que
$Y' \times_Y X \to Y'$ soit fini.
\end{lemma}

\begin{proof}
D'après le lemme \ref{lemma-quasi-finite-in-codim-1}, le morphisme $f$ est
quasi-fini en tout point situé au-dessus de $y$. Soit $\overline{y} : \Spec(k) \to Y$
un point géométrique situé au-dessus de $y$. Alors $|X_{\overline{y}}|$ est un espace
discret (Espaces décents, Lemme
\ref{decent-spaces-lemma-conditions-on-fibre-and-qf}).
L'espace $X_{\overline{y}}$ étant quasi-compact puisque $f$ est propre, nous concluons
que $|X_{\overline{y}}|$ est fini.
Nous pouvons donc appliquer Cohomologie des espaces, Lemme
\ref{spaces-cohomology-lemma-proper-finite-fibre-finite-in-neighbourhood}
pour conclure.
\end{proof}

\begin{lemma}
\label{lemma-modification-normal-iso-over-codimension-1}
Soit $S$ un schéma. Soit $X$ un espace algébrique noethérien sur $S$.
Soit $f : Y \to X$ un morphisme propre birationnel d'espaces algébriques
tel que $Y$ soit réduit.
Soit $U \subset X$ l'ouvert maximal sur lequel $f$ est un isomorphisme.
Alors $U$ contient
\begin{enumerate}
\item tout point de codimension $0$ de $X$,
\item tout $x \in |X|$ de codimension $1$ de $X$ tel que l'anneau local
de $X$ en $x$ soit normal (Propriétés des espaces, Remarque
\ref{spaces-properties-remark-list-properties-local-ring-local-etale-topology}),
et
\item tout $x \in |X|$ tel que la fibre de $|Y| \to |X|$ au-dessus de $x$ soit
finie et tel que l'anneau local de $X$ en $x$ soit normal.
\end{enumerate}
\end{lemma}

\begin{proof}
La partie (1) résulte du lemme
\ref{decent-spaces-lemma-birational-isomorphism-over-dense-open} du chapitre Espaces décents
(et du fait que les espaces algébriques noethériens $X$ et $Y$
sont quasi-séparés et donc décents).
La partie (2) résulte de la partie (3) et du lemme \ref{lemma-finite-in-codim-1}
(et du fait que les morphismes finis ont des fibres finies).
Soit $x \in |X|$ satisfaisant (3). D'après
Cohomologie des espaces, Lemme
\ref{spaces-cohomology-lemma-proper-finite-fibre-finite-in-neighbourhood}
(qui s'applique d'après le lemme
\ref{decent-spaces-lemma-conditions-on-fibre-and-qf} du chapitre Espaces décents),
nous pouvons supposer que $f$ est fini. Choisissons un schéma affine $X'$ et
un morphisme \'etale $X' \to X$ et un point $x' \in X$ s'envoyant sur $x$.
Il suffit de montrer qu'il existe un voisinage ouvert $U'$ de $x' \in X'$
tel que $Y \times_X X' \to X'$ soit un isomorphisme au-dessus de $U'$
(en effet, $U$ contient alors l'image de $U'$ dans $X$, voir Espaces, Lemme
\ref{spaces-lemma-descent-representable-transformations-property}).
Alors $Y \times_X X' \to X$ est un morphisme fini et birationnel
(Espaces décents, Lemme \ref{decent-spaces-lemma-birational-etale-localization}).
Comme un morphisme fini est affine, nous nous ramenons
au cas d'un morphisme fini birationnel de schémas affines noethériens
$Y \to X$ et d'un point $x \in X$ tel que $\mathcal{O}_{X, x}$ soit un
anneau intègre normal. Ceci est traité dans Variétés, Lemme
\ref{varieties-lemma-modification-normal-iso-over-codimension-1}.
\end{proof}






\section{Espaces algébriques intègres}
\label{section-integral-spaces}

\noindent
Nous n'avons pas encore défini la notion d'espace algébrique intègre. La
difficulté vient de ce que la propriété d'être intègre n'est pas locale pour la topologie \'etale des schémas.
Nous pourrions prendre comme définition le fait que $X$ soit réduit et que $|X|$ soit irréductible,
caractérisation donnée dans Propriétés, Lemme \ref{properties-lemma-characterize-integral}.
Dans ce cas, l'espace algébrique
décrit dans Espaces, Exemple \ref{spaces-example-infinite-product}
serait intègre, ce qui ne paraît pas souhaitable.
Pour éviter ce type de pathologie, nous supposerons de plus que $X$ est
un espace algébrique décent, bien qu'une hypothèse plus faible puisse peut-être suffire.

\begin{definition}
\label{definition-integral-algebraic-space}
Soit $S$ un schéma. On dit qu'un espace algébrique $X$ sur $S$ est
{\it intègre} s'il est réduit et décent, et si $|X|$ est irréductible.
\end{definition}

\noindent
Dans ce cas, l'espace topologique irréductible $|X|$ est sobre
(Espaces décents, Proposition \ref{decent-spaces-proposition-reasonable-sober}).
Il admet donc un unique point générique $x$. En fait, le lemme
\ref{decent-spaces-lemma-finitely-many-irreducible-components} du chapitre
Espaces décents caractérise les espaces algébriques décents
ayant un nombre fini de
composantes irréductibles. Il en résulte qu'un
espace algébrique $X$ est intègre s'il est
réduit, s'il admet un sous-schéma ouvert dense et irréductible $X'$,
de point générique $x'$, et si le morphisme $x' \to X$ est quasi-compact.

\begin{lemma}
\label{lemma-integral-algebraic-space-rational-functions}
Soit $S$ un schéma. Soit $X$ un espace algébrique intègre sur $S$.
Soit $\eta \in |X|$ le point générique de $X$.
On a des identifications canoniques
$$
R(X) = \mathcal{O}_{X, \eta}^h = \kappa(\eta)
$$
où $R(X)$ désigne l'anneau des fonctions rationnelles défini dans
Morphismes d'espaces, Définition
\ref{spaces-morphisms-definition-ring-of-rational-functions},
$\kappa(\eta)$ est le corps résiduel défini dans
Espaces décents, Définition \ref{decent-spaces-definition-residue-field},
et $\mathcal{O}_{X, \eta}^h$ est l'anneau local hensélien défini dans
Espaces décents, Définition
\ref{decent-spaces-definition-elemenary-etale-neighbourhood}.
En particulier, ces anneaux sont des corps.
\end{lemma}

\begin{proof}
Comme $X$ est un schéma dans un voisinage ouvert de $\eta$ (voir ci-dessus),
l'assertion découle immédiatement du résultat correspondant pour
les schémas ; voir Morphismes, Lemme
\ref{morphisms-lemma-integral-scheme-rational-functions}.
Nous utilisons aussi le fait que l'hensélisation d'un corps est ce corps lui-même,
ainsi que la compatibilité de nos définitions de ces objets
pour les espaces algébriques avec celles qui valent pour les schémas.
Nous omettons les détails.
\end{proof}

\noindent
Cela conduit à la définition suivante.

\begin{definition}
\label{definition-function-field}
Soit $S$ un schéma. Soit $X$ un espace algébrique intègre sur $S$.
Le {\it corps des fonctions}, ou le {\it corps des fonctions rationnelles},
de $X$ est le corps $R(X)$ du
lemme \ref{lemma-integral-algebraic-space-rational-functions}.
\end{definition}

\noindent
Nous noterons parfois ce corps $k(X)$ au lieu de $R(X)$.

\begin{lemma}
\label{lemma-integral-sections}
Soit $S$ un schéma. Soit $X$ un espace algébrique intègre sur $S$.
Alors $\Gamma(X, \mathcal{O}_X)$ est un anneau intègre.
\end{lemma}

\begin{proof}
Posons $R = \Gamma(X, \mathcal{O}_X)$. Si $f, g \in R$ sont non nuls et
$fg = 0$, alors $X = V(f) \cup V(g)$, où $V(f)$ désigne le sous-espace fermé
de $X$ défini par $f$. Comme $X$ est irréductible, on a soit
$V(f) = X$, soit $V(g) = X$. Il s'ensuit que soit $f = 0$, soit $g = 0$, d'après
Propriétés des espaces, Lemme \ref{spaces-properties-lemma-reduced-space}.
\end{proof}

\noindent
Voici un lemme sur les espaces algébriques intègres et normaux.

\begin{lemma}
\label{lemma-normal-integral-cover-by-affines}
Soit $S$ un schéma. Soit $X$ un espace algébrique intègre et normal sur $S$.
Pour tout $x \in |X|$, il existe un schéma affine intègre et normal $U$
et un morphisme \'etale $U \to X$ dont l'image contient $x$.
\end{lemma}

\begin{proof}
Choisissons un schéma affine $U$ et un morphisme \'etale $U \to X$ dont
l'image contient $x$. Soient $u_i$, $i \in I$, les points génériques des
composantes irréductibles de $U$. Alors chaque $u_i$ a pour image le point générique de $X$
(Espaces décents, Lemme \ref{decent-spaces-lemma-decent-generic-points}). D'après
notre définition d'un espace décent
(Espaces décents, Définition \ref{decent-spaces-definition-very-reasonable}),
on voit que $I$ est fini. Ainsi, $U = \Spec(A)$, où $A$ est un anneau normal
ayant un nombre fini d'idéaux premiers minimaux.
Par conséquent, $A = \prod_{i \in I} A_i$ est un produit d'anneaux intègres normaux d'après
Algèbre, Lemme \ref{algebra-lemma-characterize-reduced-ring-normal}.
Alors $U = \coprod U_i$ avec $U_i = \Spec(A_i)$ et $x$ appartient à l'image de
$U_i \to X$ pour un certain $i$. Cela démontre le lemme.
\end{proof}

\begin{lemma}
\label{lemma-normal-integral-sections}
Soit $S$ un schéma. Soit $X$ un espace algébrique intègre et normal sur $S$.
Alors $\Gamma(X, \mathcal{O}_X)$ est un anneau intègre normal.
\end{lemma}

\begin{proof}
Posons $R = \Gamma(X, \mathcal{O}_X)$. Alors $R$ est un anneau intègre d'après
Lemme \ref{lemma-integral-sections}.
Soit $f = a/b$ un élément du corps des fractions de $R$
qui est entier sur $R$.
Pour tout $U \to X$ \'etale, où $U$ est un schéma, il existe au plus un
$f_U \in \Gamma(U, \mathcal{O}_U)$ tel que $b|_U f_U = a|_U$.
En effet, $U$ est réduit et les points génériques de $U$ s'envoient sur
le point générique de $X$, ce qui implique que $b|_U$ est un
non diviseur de zéro.
Pour tout $x \in |X|$, choisissons $U \to X$ comme dans
Lemme \ref{lemma-normal-integral-cover-by-affines}.
Il existe alors un unique $f_U \in \Gamma(U, \mathcal{O}_U)$
tel que $b|_U f_U = a|_U$, puisque
$\Gamma(U, \mathcal{O}_U)$ est un anneau intègre normal d'après
Propriétés, Lemme \ref{properties-lemma-normal-integral-sections}.
Par l'unicité ci-dessus, ces $f_U$
se recollent en une section globale $f$ du faisceau
structural, c'est-à-dire en un élément de $R$.
\end{proof}

\begin{lemma}
\label{lemma-decent-irreducible-closed}
Soit $S$ un schéma. Soit $X$ un espace algébrique décent sur $S$.
On a des bijections canoniques entre les ensembles suivants :
\begin{enumerate}
\item l'ensemble des points de $X$, c'est-à-dire $|X|$,
\item l'ensemble des parties fermées irréductibles de $|X|$,
\item l'ensemble des sous-espaces fermés intègres de $X$.
\end{enumerate}
La bijection de (1) sur (2) envoie $x$ sur $\overline{\{x\}}$.
La bijection de (3) sur (2) envoie $Z$ sur $|Z|$.
\end{lemma}

\begin{proof}
L'application ainsi définie est une bijection entre (1) et (2), puisque $|X|$ est
sobre d'après
Espaces décents, Proposition \ref{decent-spaces-proposition-reasonable-sober}.
Soit $T \subset |X|$ une partie fermée irréductible. Il existe un
unique sous-espace fermé réduit $Z \subset X$ tel que
$|Z| = T$ : $Z$ est le sous-espace porté par $T$ et muni de la structure réduite
induite ; voir Propriétés des espaces, Définition
\ref{spaces-properties-definition-reduced-induced-space}.
C'est un espace algébrique intègre car il est décent,
réduit et irréductible.
\end{proof}






\section{Morphismes entre espaces algébriques intègres}
\label{section-morphisms-between-integral-spaces}

\noindent
Le lemme suivant caractérise les morphismes dominants de degré fini
entre espaces algébriques intègres.

\begin{lemma}
\label{lemma-finite-degree}
Soit $S$ un schéma. Soient $X$, $Y$ des espaces algébriques intègres sur $S$.
Soient $x \in |X|$ et $y \in |Y|$ les points génériques. Soit $f : X \to Y$
un morphisme localement de type fini. Supposons que $f$ soit dominant
(Morphismes d'espaces, Définition \ref{spaces-morphisms-definition-dominant}).
Les conditions suivantes sont équivalentes :
\begin{enumerate}
\item le degré de transcendance de $x/y$ est $0$,
\item l'extension $\kappa(x)/\kappa(y)$ (voir la démonstration) est finie,
\item il existe des ouverts affines non vides $U \subset X$ et $V \subset Y$
tels que $f(U) \subset V$ et que $f|_U : U \to V$ soit fini,
\item $f$ est quasi-fini en $x$, et
\item $x$ est le seul point de $|X|$ ayant pour image $y$.
\end{enumerate}
Si $f$ est séparé ou si $f$ est quasi-compact, alors ces conditions sont
également équivalentes à
\begin{enumerate}
\item[(6)] il existe un ouvert affine non vide $V \subset Y$ tel
que $f^{-1}(V) \to V$ soit fini.
\end{enumerate}
\end{lemma}

\begin{proof}
Un argument de topologie élémentaire montre que $f(x) = y$, puisque $f$ est dominant.
Soit $Y' \subset Y$ le lieu schématique de $Y$ et soit
$X' \subset f^{-1}(Y')$ le lieu schématique de $f^{-1}(Y')$.
D'après la discussion ci-dessus, en utilisant
Espaces décents, Proposition \ref{decent-spaces-proposition-reasonable-sober} et
Théorème \ref{decent-spaces-theorem-decent-open-dense-scheme},
on voit que $x \in |X'|$ et $y \in |Y'|$.
Alors $f|_{X'} : X' \to Y'$ est un morphisme de schémas intègres
qui est localement de type fini. On en déduit que les conditions (1), (2) et (3)
sont équivalentes d'après Morphismes, Lemme \ref{morphisms-lemma-finite-degree}.

\medskip\noindent
La condition (4) implique la condition (1) d'après
Morphismes d'espaces, Lemme \ref{spaces-morphisms-lemma-compare-tr-deg}
appliqué à $X \to Y \to Y$.
Réciproquement, la condition (3) implique la condition (4), car tout morphisme
fini est quasi-fini et $x \in U$, puisque $x$
est le point générique. Ainsi (1) -- (4) sont équivalentes.

\medskip\noindent
Supposons satisfaites les conditions équivalentes (1) -- (4). Supposons que
$x' \mapsto y$. Alors $x \leadsto x'$ est une spécialisation dans la
fibre de $|X| \to |Y|$ au-dessus de $y$. Si $x' \not = x$, alors $f$ n'est pas
quasi-fini en $x$ d'après Espaces décents, Lemme
\ref{decent-spaces-lemma-conditions-on-point-in-fibre-and-qf}.
Donc $x = x'$ et (5) est vérifiée. Réciproquement, si (5) est vérifiée, alors
(5) est vérifiée pour le morphisme de schémas $X' \to Y'$ (voir ci-dessus)
et nous pouvons utiliser
Morphismes, Lemme \ref{morphisms-lemma-finite-degree}
pour voir que (1) est vérifiée.

\medskip\noindent
Remarquons que (6) implique les conditions équivalentes (1) -- (5)
sans hypothèse supplémentaire sur $f$. Pour achever la démonstration,
il faut montrer que les conditions équivalentes (1) -- (5) impliquent (6).
Cela résulte du lemme
\ref{decent-spaces-lemma-finite-over-dense-open} du chapitre Espaces décents.
\end{proof}

\begin{definition}
\label{definition-degree}
Soit $S$ un schéma.
Soient $X$ et $Y$ des espaces algébriques intègres sur $S$.
Soit $f : X \to Y$ un morphisme localement de type fini et dominant.
Supposons que l'une des conditions équivalentes (1) -- (5) du
lemme \ref{lemma-finite-degree} soit satisfaite. Soient $x \in |X|$ et $y \in |Y|$
les points génériques. Alors l'entier positif
$$
\deg(X/Y) = [\kappa(x) : \kappa(y)]
$$
est appelé le {\it degré de $X$ sur $Y$}.
\end{definition}

\begin{lemma}
\label{lemma-degree-composition}
Soit $S$ un schéma.
Soient $X$, $Y$, $Z$ des espaces algébriques intègres sur $S$.
Soient $f : X \to Y$ et $g : Y \to Z$ des morphismes dominants et localement
de type fini. Supposons que l'une des conditions équivalentes
(1) -- (5) du lemme \ref{lemma-finite-degree} soit satisfaite pour $f$ et $g$. Alors
$$
\deg(X/Z) = \deg(X/Y) \deg(Y/Z).
$$
\end{lemma}

\begin{proof}
Cela découle de la multiplicativité des degrés dans les tours
d'extensions finies de corps ; voir
Corps, Lemme \ref{fields-lemma-multiplicativity-degrees}.
\end{proof}










\section{Diviseurs de Weil}
\label{section-Weil-divisors}

\noindent
Cette section est l'analogue de la section
\ref{divisors-section-Weil-divisors} de Diviseurs.

\medskip\noindent
Nous allons introduire les diviseurs de Weil et l'équivalence rationnelle des
diviseurs de Weil pour les espaces algébriques intègres localement noethériens.
Comme nous ne supposons pas que nos espaces algébriques sont quasi-compacts, nous devons
être un peu prudents lors de la définition des diviseurs de Weil. Nous devons permettre
des sommes infinies de diviseurs premiers, car une fonction rationnelle peut avoir,
par exemple, une infinité de pôles. Dans le cas quasi-compact, nos diviseurs
de Weil sont des sommes finies comme d'habitude. Voici un lemme fondamental que nous
utiliserons souvent pour montrer que des familles de sous-espaces fermés sont localement finies.

\begin{lemma}
\label{lemma-components-locally-finite}
Soit $S$ un schéma et soit $X$ un espace algébrique localement noethérien
sur $S$. Si $T \subset |X|$ est une partie fermée,
alors la famille des composantes irréductibles de $T$ est localement finie.
\end{lemma}

\begin{proof}
L'espace topologique $|X|$ est localement noethérien
(Propriétés des espaces, lemme \ref{spaces-properties-lemma-Noetherian-topology}).
Un espace topologique noethérien possède un nombre fini de
composantes irréductibles et tout sous-espace d'un espace noethérien est noethérien
(Topologie, lemme \ref{topology-lemma-Noetherian}).
Le lemme découle donc de la définition d'une famille localement finie
(Topologie, définition \ref{topology-definition-locally-finite}).
\end{proof}

\noindent
Soit $S$ un schéma. Soit $X$ un espace algébrique décent sur $S$.
Soit $Z$ un sous-espace fermé intègre de $X$ et soit
$\xi \in |Z|$ son point générique. Alors la codimension de
$|Z|$ dans $|X|$ est égale à la dimension de l'anneau local
de $X$ en $\xi$, d'après
le lemme d'Espaces décents \ref{decent-spaces-lemma-codimension-local-ring}.
Rappelons que nous indiquons également cela en disant que
{\it $\xi$ est un point de codimension $1$ sur $X$}, voir la
définition de Propriétés des espaces
\ref{spaces-properties-definition-dimension-local-ring}.

\begin{definition}
\label{definition-Weil-divisor}
Soit $S$ un schéma.
Soit $X$ un espace algébrique intègre localement noethérien sur $S$.
\begin{enumerate}
\item Un {\it diviseur premier} est un sous-espace fermé intègre $Z \subset X$
de codimension $1$, c'est-à-dire que le point générique de $|Z|$ est un point
de codimension $1$ sur $X$.
\item Un {\it diviseur de Weil} est une somme formelle $D = \sum n_Z Z$ où
la somme porte sur les diviseurs premiers de $X$ et la famille
$\{|Z| : n_Z \not = 0\}$ est localement finie dans $|X|$
(Topologie, définition \ref{topology-definition-locally-finite}).
\end{enumerate}
Le groupe de tous les diviseurs de Weil sur $X$ est noté $\text{Div}(X)$.
\end{definition}

\noindent
Notre prochaine tâche consiste à définir le diviseur de Weil associé à une fonction
rationnelle. Pour cela, nous devons définir l'ordre d'annulation d'une
fonction rationnelle sur un espace algébrique intègre localement noethérien $X$
le long d'un diviseur premier $Z$. Soit $\xi \in |Z|$ le point générique.
Nous rencontrons ici le problème que l'anneau local $\mathcal{O}_{X, \xi}$
n'existe pas et que l'anneau local hensélien $\mathcal{O}_{X, \xi}^h$
peut ne pas être un anneau intègre, voir l'exemple \ref{example-not-a-domain}.
Pour contourner cela, nous utilisons le lemme suivant.

\begin{lemma}
\label{lemma-order-vanishing}
Soit $S$ un schéma. Soit $X$ un espace algébrique intègre localement noethérien
sur $S$. Soit $Z \subset X$ un diviseur premier et soit $\xi \in |Z|$
le point générique. Alors l'anneau local hensélien $\mathcal{O}_{X, \xi}^h$
est un anneau local noethérien réduit de dimension $1$ et
il existe une application canonique injective
$$
R(X) \longrightarrow Q(\mathcal{O}_{X, \xi}^h)
$$
du corps des fonctions $R(X)$ de $X$ dans l'anneau total des fractions.
\end{lemma}

\begin{proof}
Nous utiliserons les résultats d'Espaces décents, section
\ref{decent-spaces-section-residue-fields-henselian-local-rings}.
Soit $(U, u) \to (X, \xi)$ un voisinage \'etale élémentaire.
Observons que $U$ est localement noethérien et réduit.
Ainsi, $\mathcal{O}_{U, u}$ est un anneau de dimension $1$
(par notre définition des diviseurs premiers),
réduit et noethérien.
Après avoir remplacé $U$ par un voisinage ouvert affine de $u$
nous pouvons supposer que $U$ est noethérien et affine.
Après avoir remplacé $U$ par un ouvert plus petit, nous pouvons supposer que toute
composante irréductible de $U$ passe par $u$.
Comme $U \to X$ est ouvert et que $X$ est irréductible, $U \to X$ est dominant.
Nous obtenons donc une application d'anneaux $R(X) \to R(U)$ en composant des applications rationnelles, voir
Morphismes d'espaces, section \ref{spaces-morphisms-section-rational-maps}.
Comme $R(X)$ est un corps, cette application est injective.
Par notre choix de $U$, nous voyons que $R(U)$ est l'anneau total des fractions
$Q(\mathcal{O}_{U, u})$, voir
Morphismes, lemme \ref{morphisms-lemma-integral-scheme-rational-functions}
et
Algèbre, lemme \ref{algebra-lemma-total-ring-fractions-no-embedded-points}.

\medskip\noindent
À ce stade, nous avons démontré toutes les assertions du lemme avec
$\mathcal{O}_{U, u}$ à la place de $\mathcal{O}_{X, \xi}^h$.
Cependant, $\mathcal{O}_{X, \xi}^h$ est l'hensélisation de
$\mathcal{O}_{U, u}$. Ainsi, $\mathcal{O}_{X, \xi}^h$ est un anneau
réduit noethérien de dimension $1$, voir
Compléments d'algèbre, lemmes
\ref{more-algebra-lemma-henselization-reduced},
\ref{more-algebra-lemma-henselization-dimension}, et
\ref{more-algebra-lemma-henselization-noetherian}.
Comme $\mathcal{O}_{U, u} \to \mathcal{O}_{X, \xi}^h$ est
fidèlement plat d'après Compléments d'algèbre, lemme
\ref{more-algebra-lemma-dumb-properties-henselization}
elle envoie les non diviseurs de zéro sur des non diviseurs de zéro.
Nous obtenons donc une application canonique
$Q(\mathcal{O}_{U, u}) \to Q(\mathcal{O}_{X, \xi}^h)$
et nous obtenons notre application.
Nous omettons la vérification que l'application est indépendante du
choix de $(U, u) \to (X, x)$ ; une approche légèrement meilleure
consisterait d'abord à observer que
$\colim Q(\mathcal{O}_{U, u}) = Q(\mathcal{O}_{X, \xi}^h)$.
\end{proof}

\begin{definition}
\label{definition-order-vanishing}
Soit $S$ un schéma. Soit $X$ un espace algébrique intègre localement noethérien
sur $S$. Soit $f \in R(X)^*$. Pour tout diviseur premier
$Z \subset X$, nous définissons l'{\it ordre d'annulation de $f$ le long de $Z$}
comme l'entier
$$
\text{ord}_Z(f) =
\text{longueur}_{\mathcal{O}_{X, \xi}^h}
(\mathcal{O}_{X, \xi}^h/a \mathcal{O}_{X, \xi}^h) -
\text{longueur}_{\mathcal{O}_{X, \xi}^h}
(\mathcal{O}_{X, \xi}^h/b \mathcal{O}_{X, \xi}^h)
$$
où $a, b \in \mathcal{O}_{X, \xi}^h$ sont des non diviseurs de zéro
tels que l'image de $f$ dans $Q(\mathcal{O}_{X, \xi}^h)$
(lemme \ref{lemma-order-vanishing}) soit égale à $a/b$.
Ceci est bien défini d'après
Algèbre, lemme \ref{algebra-lemma-ord-additive}.
\end{definition}

\noindent
Si $\mathcal{O}_{X, \xi}^h$ est un anneau intègre, nous obtenons
$$
\text{ord}_Z(f) = \text{ord}_{\mathcal{O}_{X, \xi}^h}(f)
$$
où le membre de droite est la notion de
la définition d'Algèbre \ref{algebra-definition-ord}.
Notons que pour $f, g \in R(X)^*$ nous avons
$$
\text{ord}_Z(fg) = \text{ord}_Z(f) + \text{ord}_Z(g).
$$
Bien sûr, il peut arriver que $\text{ord}_Z(f) < 0$.
Dans ce cas, nous disons que $f$ possède un {\it pôle} le long de $Z$
et que $-\text{ord}_Z(f) > 0$ est l'
{\it ordre du pôle de $f$ le long de $Z$}. Il est important de remarquer que
la condition $\text{ord}_Z(f) \geq 0$ n'est {\bf pas} équivalente
à la condition $f \in \mathcal{O}_{X, \xi}^h$, sauf si l'anneau local
$\mathcal{O}_{X, \xi}$ est un anneau de valuation discrète.

\begin{lemma}
\label{lemma-order-vanishing-agrees}
Soit $S$ un schéma. Soit $X$ un espace algébrique intègre localement noethérien
sur $S$. Soit $f \in R(X)^*$. Si le diviseur premier
$Z \subset X$ rencontre le lieu schématique de $X$, alors l'ordre
d'annulation $\text{ord}_Z(f)$ de la définition \ref{definition-order-vanishing}
coïncide avec l'ordre d'annulation de
Diviseurs, définition \ref{divisors-definition-order-vanishing}.
\end{lemma}

\begin{proof}
Après avoir rétréci $X$, nous pouvons supposer que $X$ est un schéma intègre noethérien.
Si $\xi \in Z$ désigne le point générique, alors
$\mathcal{O}_{X, \xi}^h$ est l'hensélisation de $\mathcal{O}_{X, \xi}$
(Espaces décents, lemme \ref{decent-spaces-lemma-describe-henselian-local-ring}).
Pour démontrer le lemme, il suffit et il est nécessaire de montrer que
$$
\text{longueur}_{\mathcal{O}_{X, \xi}}
(\mathcal{O}_{X, \xi}/a \mathcal{O}_{X, \xi}) =
\text{longueur}_{\mathcal{O}_{X, \xi}^h}
(\mathcal{O}_{X, \xi}^h/a \mathcal{O}_{X, \xi}^h)
$$
Ceci découle immédiatement de
Algèbre, lemme \ref{algebra-lemma-pullback-module}
(et du fait que $\mathcal{O}_{X, \xi} \to \mathcal{O}_{X, \xi}^h$
est un homomorphisme local plat d'anneaux locaux noethériens).
\end{proof}

\begin{lemma}
\label{lemma-divisor-locally-finite}
Soit $S$ un schéma. Soit $X$ un espace algébrique intègre localement noethérien
sur $S$. Soit $f \in R(X)^*$. Alors les familles
$$
\{Z \subset X \mid Z\text{ est un diviseur premier de point générique }\xi
\text{ et }f\text{ n'appartient pas à }\mathcal{O}_{X, \xi}\}
$$
et
$$
\{Z \subset X \mid Z \text{ est un diviseur premier et }\text{ord}_Z(f) \not = 0\}
$$
sont localement finies dans $X$.
\end{lemma}

\begin{proof}
Il existe un sous-espace ouvert non vide $U \subset X$
tel que $f$ corresponde à une section de $\Gamma(U, \mathcal{O}_X^*)$.
Par conséquent, les diviseurs premiers qui peuvent apparaître dans les ensembles du lemme
correspondent tous à des composantes irréductibles de $|X| \setminus |U|$.
Le lemme \ref{lemma-components-locally-finite} donne donc le résultat voulu.
\end{proof}

\noindent
Ce lemme nous permet de donner la définition suivante.

\begin{definition}
\label{definition-principal-divisor}
Soit $S$ un schéma.
Soit $X$ un espace algébrique intègre localement noethérien sur $S$.
Soit $f \in R(X)^*$.
Le {\it diviseur de Weil principal associé à $f$} est le diviseur de Weil
$$
\text{div}(f) = \text{div}_X(f) = \sum \text{ord}_Z(f) [Z]
$$
où la somme porte sur les diviseurs premiers et où $\text{ord}_Z(f)$ est comme dans
la définition \ref{definition-order-vanishing}. Ceci a un sens
d'après le lemme \ref{lemma-divisor-locally-finite}.
\end{definition}

\begin{lemma}
\label{lemma-div-additive}
Soit $S$ un schéma.
Soit $X$ un espace algébrique intègre localement noethérien sur $S$.
Soient $f, g \in R(X)^*$. Alors
$$
\text{div}_X(fg) = \text{div}_X(f) + \text{div}_X(g)
$$
en tant que diviseurs de Weil sur $X$.
\end{lemma}

\begin{proof}
Ceci est clair par additivité des fonctions $\text{ord}$.
\end{proof}

\noindent
Nous voyons d'après le lemme ci-dessus que la famille des diviseurs de Weil principaux
forme un sous-groupe du groupe de tous les diviseurs de Weil. Cela conduit à la
définition suivante.

\begin{definition}
\label{definition-class-group}
Soit $S$ un schéma.
Soit $X$ un espace algébrique intègre localement noethérien sur $S$. Le
{\it groupe des classes de diviseurs de Weil} de $X$ est le quotient du
groupe des diviseurs de Weil par le sous-groupe des diviseurs de Weil principaux.
Notation : $\text{Cl}(X)$.
\end{definition}

\noindent
Par construction, nous obtenons un complexe exact
\begin{equation}
\label{equation-Weil-divisor-class}
R(X)^* \xrightarrow{\text{div}} \text{Div}(X) \to \text{Cl}(X) \to 0
\end{equation}
que nous pouvons considérer comme une présentation de $\text{Cl}(X)$. Notre prochaine tâche
est de relier le groupe des classes de diviseurs de Weil au groupe de Picard.

\begin{example}
\label{example-Z-mod-2}
Cet exemple prolonge celui de Morphismes d'espaces,
exemple \ref{spaces-morphisms-example-universal-homeomorphism}.
Considérons l'espace algébrique
$X = \mathbf{A}^1_k/\{t \sim -t \mid t \not = 0\}$.
C'est un espace algébrique lisse sur le corps $k$.
Il existe un homéomorphisme universel
$$
X \longrightarrow \mathbf{A}^1_k = \Spec(k[t])
$$
qui est un isomorphisme au-dessus de $\mathbf{A}^1_k \setminus \{0\}$.
Nous en concluons que $X$ est noethérien et intègre.
Comme $\dim(X) = 1$, les diviseurs premiers de $X$ sont
les points fermés de $X$.
Considérons l'unique point fermé $x \in |X|$ situé au-dessus de $0 \in \mathbf{A}^1_k$.
Comme $X \setminus \{x\}$ s'envoie isomorphiquement sur $\mathbf{A}^1 \setminus \{0\}$,
nous voyons que les classes dans $\text{Cl}(X)$ des points fermés distincts
de $x$ sont nulles. Cependant, le diviseur de $t$ sur $X$ est $2[x]$.
Nous concluons que $\text{Cl}(X) = \mathbf{Z}/2\mathbf{Z}$.
\end{example}

\begin{example}
\label{example-not-a-domain}
Soit $k$ un corps. Soit
$$
U = \Spec(k[x, y]/(xy))
$$
qui est l'union des axes de coordonnées dans $\mathbf{A}^2_k$.
Notons $\Delta : U \to U \times_k U$ la diagonale et
$\Delta' : U \to U \times_k U$ l'application $u \mapsto (u, \sigma(u))$
où $\sigma : U \to U$, $(x, y) \mapsto (y, x)$
est l'automorphisme qui échange les axes de coordonnées. Posons
$$
R = \Delta(U) \amalg \Delta'(U \setminus \{0_U\})
$$
où $0_U \in U$ est l'origine. Il est facile de voir que
$R$ est une relation d'équivalence \'etale sur $U$. Le quotient $X = U/R$
est un espace algébrique. Le morphisme $U \to \mathbf{A}^1_k$,
$(x, y) \mapsto x + y$ est invariant par $R$ et définit donc
un morphisme
$$
X \longrightarrow \mathbf{A}^1_k
$$
Ce morphisme est un homéomorphisme universel et un isomorphisme au-dessus de
$\mathbf{A}^1_k \setminus \{0\}$. Il s'ensuit que $X$ est intègre
et noethérien. Exactement comme dans l'exemple \ref{example-Z-mod-2}
le lecteur montre que $\text{Cl}(X) = \mathbf{Z}/2\mathbf{Z}$
avec pour générateur la classe correspondant à l'unique point fermé $x \in |X|$
situé au-dessus de $0 \in \mathbf{A}^1_k$. Cependant, dans ce cas, l'anneau local hensélien
de $X$ en $x$ n'est pas un anneau intègre, car c'est l'hensélisation
de $\mathcal{O}_{U, 0_U}$.
\end{example}














\section{Classe de diviseur de Weil associée à un module inversible}
\label{section-c1}

\noindent
Dans cette section, nous suivons exactement la même progression que dans la
section \ref{section-Weil-divisors} pour définir une application canonique
$\Pic(X) \to \text{Cl}(X)$
sur un espace algébrique intègre localement noethérien.

\medskip\noindent
Soit $S$ un schéma. Soit $X$ un espace algébrique intègre localement noethérien
sur $S$. Soit $\mathcal{L}$ un $\mathcal{O}_X$-module inversible.
D'après Diviseurs sur les espaces, lemme
\ref{spaces-divisors-lemma-regular-meromorphic-section-exists}
il existe une section méromorphe régulière
$s \in \Gamma(X, \mathcal{K}_X(\mathcal{L}))$.
En fait, d'après Diviseurs sur les espaces, lemme
\ref{spaces-divisors-lemma-compute-meromorphic}
c'est la même chose qu'un élément non nul de
$\mathcal{L}_\eta$, où $\eta \in |X|$ est le point générique.
Le même lemme nous dit que si $\mathcal{L} = \mathcal{O}_X$,
alors $s$ est la même chose qu'une fonction rationnelle non nulle sur $X$
(ainsi ce que nous ferons ci-dessous coïncide avec la construction de la
section \ref{section-Weil-divisors}).

\medskip\noindent
Soit $Z \subset X$ un diviseur premier et soit $\xi \in |Z|$ le point générique.
Nous allons définir l'ordre d'annulation de $s$ le long de $Z$.
Considérons le morphisme canonique
$$
c_\xi : \Spec(\mathcal{O}_{X, \xi}^h) \longrightarrow X
$$
dont la source est le spectre de l'anneau local hensélien de $X$ en $\xi$
(Espaces décents, définition \ref{decent-spaces-definition-henselian-local-ring}).
Le tiré en arrière $\mathcal{L}_\xi = c_\xi^*\mathcal{L}$
est un module inversible, donc trivial ; choisissons un générateur $s_\xi$ de
$\mathcal{L}_\xi$. Comme $c_\xi$ est plat, les tirés en arrière des fonctions méromorphes
et des sections (régulières) sont définis pour $c_\xi$, voir
Diviseurs sur les espaces, définition
\ref{spaces-divisors-definition-pullback-meromorphic-sections} et
les lemmes \ref{spaces-divisors-lemma-pullback-meromorphic-sections-defined} et
\ref{spaces-divisors-lemma-meromorphic-sections-pullback}.
Nous obtenons ainsi
$$
c_\xi^*(s) = f s_\xi
$$
pour un élément non diviseur de zéro $f \in Q(\mathcal{O}_{X, \xi}^h)$.
Nous utilisons ici le lemme de Diviseurs \ref{divisors-lemma-locally-Noetherian-K}
pour identifier l'espace des sections méromorphes de
$\mathcal{L}_\xi \cong \mathcal{O}_{\Spec(\mathcal{O}_{X, \xi}^h)}$
en termes de l'anneau total des fractions de $\mathcal{O}_{X, \xi}^h$.
Convenons de noter cet élément
$$
s/s_\xi = f \in  Q(\mathcal{O}_{X, \xi}^h)
$$
Observons que $f = s/s_\xi$
est remplacé par $uf$ où $u \in \mathcal{O}_{X, \xi}^h$
est une unité lorsque nous changeons le choix de $s_\xi$.

\begin{definition}
\label{definition-order-vanishing-meromorphic}
Soit $S$ un schéma. Soit $X$ un espace algébrique intègre localement noethérien
sur $S$. Soit $\mathcal{L}$ un
$\mathcal{O}_X$-module inversible.
Soit $s \in \Gamma(X, \mathcal{K}_X(\mathcal{L}))$
une section méromorphe régulière de $\mathcal{L}$.
Pour tout diviseur premier $Z \subset X$ dont le point générique est $\xi \in |Z|$
nous définissons l'
{\it ordre d'annulation de $s$ le long de $Z$}
comme l'entier
$$
\text{ord}_{Z, \mathcal{L}}(s) =
\text{longueur}_{\mathcal{O}_{X, \xi}^h}
(\mathcal{O}_{X, \xi}^h/a \mathcal{O}_{X, \xi}^h) -
\text{longueur}_{\mathcal{O}_{X, \xi}^h}
(\mathcal{O}_{X, \xi}^h/b \mathcal{O}_{X, \xi}^h)
$$
où $a, b \in \mathcal{O}_{X, \xi}^h$ sont des non diviseurs de zéro
tels que l'élément $s/s_\xi$ de $Q(\mathcal{O}_{X, \xi}^h)$
construit ci-dessus soit égal à $a/b$.
Ceci est bien défini d'après ce qui précède et
le lemme d'Algèbre \ref{algebra-lemma-ord-additive}.
\end{definition}

\noindent
Comme expliqué ci-dessus, une section méromorphe régulière $s$ de $\mathcal{O}_X$
peut s'écrire $s = f \cdot 1$ où $f$
est une fonction rationnelle non nulle sur $X$ et nous avons
$\text{ord}_Z(f) = \text{ord}_{Z, \mathcal{O}_X}(s)$.
Comme dans le cas des diviseurs principaux, nous avons le lemme suivant.

\begin{lemma}
\label{lemma-divisor-meromorphic-locally-finite}
Soit $S$ un schéma. Soit $X$ un espace algébrique intègre localement noethérien
sur $S$. Soit $\mathcal{L}$ un $\mathcal{O}_X$-module inversible.
Soit $s \in \mathcal{K}_X(\mathcal{L})$ une
section méromorphe régulière (c'est-à-dire non nulle) de $\mathcal{L}$. Alors les ensembles
$$
\{Z \subset X \mid Z \text{ est un diviseur premier de point générique }\xi
\text{ et }s\text{ n'appartient pas à }\mathcal{L}_\xi\}
$$
et
$$
\{Z \subset X \mid Z \text{ est un diviseur premier et }
\text{ord}_{Z, \mathcal{L}}(s) \not = 0\}
$$
sont localement finis dans $X$.
\end{lemma}

\begin{proof}
Il existe un sous-espace ouvert non vide $U \subset X$ tel que $s$
corresponde à une section de $\Gamma(U, \mathcal{L})$ qui engendre
$\mathcal{L}$ sur $U$. Par conséquent, les diviseurs premiers qui peuvent apparaître
dans les ensembles du lemme correspondent tous à des composantes irréductibles de
$|X| \setminus |U|$. Le lemme \ref{lemma-components-locally-finite}
donne le résultat voulu.
\end{proof}

\begin{lemma}
\label{lemma-divisor-meromorphic-well-defined}
Soit $S$ un schéma. Soit $X$ un espace algébrique intègre localement noethérien
sur $S$. Soit $\mathcal{L}$ un $\mathcal{O}_X$-module inversible.
Soient $s, s' \in \mathcal{K}_X(\mathcal{L})$ des sections méromorphes
non nulles de $\mathcal{L}$. Alors $f = s/s'$
est un élément de $R(X)^*$ et nous avons
$$
\sum \text{ord}_{Z, \mathcal{L}}(s)[Z]
=
\sum \text{ord}_{Z, \mathcal{L}}(s')[Z]
+
\text{div}(f)
$$
en tant que diviseurs de Weil.
\end{lemma}

\begin{proof}
C'est clair d'après les définitions.
Notons que le lemme \ref{lemma-divisor-meromorphic-locally-finite}
garantit que les sommes sont bien des diviseurs de Weil.
\end{proof}

\begin{definition}
\label{definition-divisor-invertible-sheaf}
Soit $S$ un schéma. Soit $X$ un espace algébrique intègre localement noethérien
sur $S$. Soit $\mathcal{L}$ un $\mathcal{O}_X$-module inversible.
\begin{enumerate}
\item Pour toute section méromorphe non nulle $s$ de $\mathcal{L}$
nous définissons le {\it diviseur de Weil associé à $s$} comme
$$
\text{div}_\mathcal{L}(s) =
\sum \text{ord}_{Z, \mathcal{L}}(s) [Z] \in \text{Div}(X)
$$
où la somme porte sur les diviseurs premiers. Ceci est bien défini d'après
le lemme \ref{lemma-divisor-meromorphic-locally-finite}.
\item Nous définissons la {\it classe de diviseur de Weil associée à $\mathcal{L}$}
comme l'image de $\text{div}_\mathcal{L}(s)$ dans $\text{Cl}(X)$
où $s$ est une quelconque section méromorphe non nulle de $\mathcal{L}$ sur $X$.
Ceci est bien défini d'après
le lemme \ref{lemma-divisor-meromorphic-well-defined}.
\end{enumerate}
\end{definition}

\noindent
Comme prévu, cette construction est additive par rapport au module inversible.

\begin{lemma}
\label{lemma-c1-additive}
Soit $S$ un schéma. Soit $X$ un espace algébrique intègre localement noethérien
sur $S$. Soient $\mathcal{L}$ et $\mathcal{N}$ deux $\mathcal{O}_X$-modules inversibles.
Soit $s$, resp.\ $t$, une section méromorphe non nulle de $\mathcal{L}$,
resp.\ de $\mathcal{N}$. Alors $st$ est une section méromorphe non nulle de
$\mathcal{L} \otimes_{\mathcal{O}_X} \mathcal{N}$ et
$$
\text{div}_{\mathcal{L} \otimes \mathcal{N}}(st)
=
\text{div}_\mathcal{L}(s) + \text{div}_\mathcal{N}(t)
$$
dans $\text{Div}(X)$. En particulier, la classe de diviseur de Weil de
$\mathcal{L} \otimes_{\mathcal{O}_X} \mathcal{N}$ est la somme
des classes de diviseurs de Weil de $\mathcal{L}$ et de $\mathcal{N}$.
\end{lemma}

\begin{proof}
Soit $s$, resp.\ $t$, une section méromorphe non nulle de
$\mathcal{L}$, resp.\ de $\mathcal{N}$. Alors $st$ est une section
méromorphe non nulle de $\mathcal{L} \otimes \mathcal{N}$.
Soit $Z \subset X$ un diviseur premier. Soit $\xi \in |Z|$ son point
générique. Choisissons des générateurs $s_\xi \in \mathcal{L}_\xi$ et
$t_\xi \in \mathcal{N}_\xi$, avec les notations introduites plus haut
dans cette section. Alors $s_\xi \otimes t_\xi$ est un générateur
de $(\mathcal{L} \otimes \mathcal{N})_\xi$.
Ainsi, $st/(s_\xi t_\xi) = (s/s_\xi)(t/t_\xi)$ dans
$Q(\mathcal{O}_{X, \xi}^h)$. En appliquant l'additivité du
lemme d'Algèbre \ref{algebra-lemma-ord-additive},
nous concluons que
$$
\text{div}_{\mathcal{L} \otimes \mathcal{N}, Z}(st)
=
\text{div}_{\mathcal{L}, Z}(s) + \text{div}_{\mathcal{N}, Z}(t)
$$
Nous omettons certains détails.
\end{proof}

\noindent
Soit $S$ un schéma. Soit $X$ un espace algébrique intègre localement noethérien
sur $S$. Les constructions et les lemmes ci-dessus donnent un homomorphisme
de groupes abéliens
\begin{equation}
\label{equation-c1}
\Pic(X) \longrightarrow \text{Cl}(X)
\end{equation}
qui associe à tout module inversible sa classe de diviseur de Weil.

\begin{lemma}
\label{lemma-normal-c1-injective}
Soit $S$ un schéma. Soit $X$ un espace algébrique intègre localement noethérien
sur $S$. Si $X$ est normal, alors l'application (\ref{equation-c1})
$\Pic(X) \to \text{Cl}(X)$ est injective.
\end{lemma}

\begin{proof}
Soit $\mathcal{L}$ un $\mathcal{O}_X$-module inversible dont la
classe de diviseur de Weil associée est triviale. Soit $s$ une section
méromorphe régulière de $\mathcal{L}$. L'hypothèse signifie que
$\text{div}_\mathcal{L}(s) = \text{div}(f)$ pour un certain
$f \in R(X)^*$. Alors $t = f^{-1}s$ est une section
méromorphe régulière de $\mathcal{L}$ telle que
$\text{div}_\mathcal{L}(t) = 0$, voir le
lemme \ref{lemma-divisor-meromorphic-well-defined}.
Nous affirmons que $t$ définit une trivialisation de $\mathcal{L}$.
Cette assertion achève la démonstration du lemme.
La démonstration qui suit est un peu maladroite, car nous ne
disposons pas encore d'une théorie très développée ; nous suggérons
au lecteur de la passer.

\medskip\noindent
Nous pouvons vérifier cette assertion localement pour la topologie \'etale. Soit
$U \in X_\etale$ affine tel que $\mathcal{L}|_U$ soit trivial. Soit
$s_U \in \Gamma(U, \mathcal{L}|_U)$ une trivialisation. D'après le
lemme de Propriétés \ref{properties-lemma-normal-locally-finite-nr-irreducibles},
nous pouvons aussi supposer $U$ intègre. Écrivons $U = \Spec(A)$, où $A$ est
un anneau intègre normal noethérien de corps des fractions $K$.
Nous pouvons écrire $t|_U = f s_U$ pour un élément $f$ de $K$, voir par exemple le
lemme de Diviseurs sur les espaces \ref{spaces-divisors-lemma-meromorphic-quasi-coherent}.
Soit $\mathfrak p \subset A$ un idéal premier de hauteur un
correspondant à un point de codimension $1$, $u \in U$, dont l'image
est un point de codimension $1$, $\xi \in |X|$. Choisissons une trivialisation
$s_\xi$ de $\mathcal{L}_\xi$ comme au début de cette section.
Choisissons un point géométrique $\overline{u}$ de $U$ situé au-dessus de $u$.
Alors
$$
(\mathcal{O}_{X, \xi}^h)^{sh} =
\mathcal{O}_{X, \overline{u}} =
\mathcal{O}_{U, u}^{sh} = (A_\mathfrak p)^{sh}
$$
voir les lemmes d'Espaces décents
\ref{decent-spaces-lemma-henselian-local-ring-strict}
et le lemme de
Propriétés des espaces
\ref{spaces-properties-lemma-describe-etale-local-ring}.
La normalité de $X$ montre que tous ces anneaux sont des
anneaux de valuation discrète. Les trivialisations $s_U$ et $s_\xi$
diffèrent par une unité comme sections du tiré en arrière de $\mathcal{L}$ sur
$\Spec(\mathcal{O}_{X, \overline{u}})$.
Écrivons $t = f_\xi s_\xi$ avec $f_\xi \in Q(\mathcal{O}_{X, \xi}^h)$.
Nous en concluons que $f_\xi$ et $f$ diffèrent par une unité
dans $Q(\mathcal{O}_{X, \overline{u}})$.
Si $Z \subset X$ désigne le diviseur premier correspondant à $\xi$
(lemme \ref{lemma-decent-irreducible-closed}), alors
$0 = \text{ord}_{Z, \mathcal{L}}(t) =
\text{ord}_{\mathcal{O}_{X, \xi}^h}(f_\xi)$
et, comme $\mathcal{O}_{X, \xi}^h$ est un anneau de valuation discrète,
nous voyons que $f_\xi$ est une unité.
Ainsi, $f$ est une unité dans $\mathcal{O}_{X, \overline{u}}$
et, en particulier, $f \in A_\mathfrak p$.
Cela implique $f \in A$ d'après le lemme d'Algèbre
\ref{algebra-lemma-normal-domain-intersection-localizations-height-1}.
Nous concluons que $t \in \Gamma(X, \mathcal{L})$.
En répétant le raisonnement avec $t^{-1}$ considéré comme section méromorphe
de $\mathcal{L}^{\otimes -1}$, on achève la démonstration.
\end{proof}

















\section{Modifications et altérations}
\label{section-modifications-alterations}

\noindent
Notre notion d'espace algébrique intègre permet de définir une modification
comme suit.

\begin{definition}
\label{definition-modification}
Soit $S$ un schéma. Soit $X$ un espace algébrique intègre sur $S$. Une
{\it modification de $X$} est un morphisme propre birationnel
$f : X' \to X$ d'espaces algébriques sur $S$, où $X'$ est intègre.
\end{definition}

\noindent
Pour les morphismes birationnels d'espaces algébriques, voir la
définition d'Espaces décents \ref{decent-spaces-definition-birational}.

\begin{lemma}
\label{lemma-modification-iso-over-open}
Soit $f : X' \to X$ une modification comme dans la
définition \ref{definition-modification}.
Il existe un ouvert non vide $U \subset X$ tel que $f^{-1}(U) \to U$
soit un isomorphisme.
\end{lemma}

\begin{proof}
D'après le
lemme \ref{lemma-finite-degree}, il existe un ouvert non vide $U \subset X$ tel
que $f^{-1}(U) \to U$ soit fini. Par platitude générique
(Morphismes d'espaces, proposition
\ref{spaces-morphisms-proposition-generic-flatness-reduced}),
nous pouvons supposer $f^{-1}(U) \to U$ plat et de présentation finie.
Ainsi, $f^{-1}(U) \to U$ est fini localement libre
(Morphismes d'espaces, lemme \ref{spaces-morphisms-lemma-finite-flat}).
Comme $f$ est birationnel, le degré de $X'$ sur $X$ est $1$.
Par conséquent, $f^{-1}(U) \to U$ est fini localement libre de degré $1$,
autrement dit, c'est un isomorphisme.
\end{proof}

\begin{definition}
\label{definition-alteration}
Soit $S$ un schéma. Soit $X$ un espace algébrique intègre sur $S$.
Une {\it altération de $X$} est un morphisme propre dominant $f : Y \to X$
d'espaces algébriques sur $S$, avec $Y$ intègre, tel que $f^{-1}(U) \to U$
soit fini pour un certain ouvert non vide $U \subset X$.
\end{definition}

\noindent
Si $f : Y \to X$ est un morphisme dominant et propre entre espaces
algébriques intègres, alors c'est une altération dès que l'extension induite
des corps résiduels aux points génériques est finie. Voici l'énoncé
précis.

\begin{lemma}
\label{lemma-alteration-generically-finite}
Soit $S$ un schéma. Soit $f : X \to Y$ un morphisme propre dominant
d'espaces algébriques intègres sur $S$. Alors $f$ est une altération
si et seulement si l'une des conditions équivalentes (1) -- (6) du
lemme \ref{lemma-finite-degree} est satisfaite.
\end{lemma}

\begin{proof}
C'est une conséquence immédiate du lemme cité dans l'énoncé.
\end{proof}

\begin{lemma}
\label{lemma-alteration-contained-in}
Soit $S$ un schéma. Soit $f : X \to Y$ un morphisme propre surjectif
d'espaces algébriques sur $S$. Supposons $Y$ intègre. Alors
il existe un sous-espace fermé intègre $X' \subset X$ tel que
$f' = f|_{X'} : X' \to Y$ soit une altération.
\end{lemma}

\begin{proof}
Soit $V \subset Y$ un ouvert affine non vide
(Espaces décents, théorème \ref{decent-spaces-theorem-decent-open-dense-scheme}).
Soit $\eta \in V$ le point générique. Alors
$X_\eta$ est un espace algébrique propre non vide sur $\eta$.
Choisissons un point fermé $x \in |X_\eta|$
(il en existe, car $|X_\eta|$ est un espace topologique quasi-compact
et sobre ; voir Espaces décents, proposition
\ref{decent-spaces-proposition-reasonable-sober},
et Topologie, lemme \ref{topology-lemma-quasi-compact-closed-point}).
Soit $X'$ le sous-espace fermé, muni de la structure réduite induite, porté par
$\overline{\{x\}} \subset |X|$ (Propriétés des espaces, définition
\ref{spaces-properties-definition-reduced-induced-space}).
Alors $f' : X' \to Y$ est surjectif, puisque son image contient $\eta$.
De plus, $f'$ est propre comme composé d'une immersion fermée
et d'un morphisme propre. Enfin, la fibre $X'_\eta$ ne possède qu'un
seul point ; pour le voir, appliquons le
lemme d'Espaces décents \ref{decent-spaces-lemma-topology-fibre}
à $X \to Y$ et à $X' \to Y$, au point $\eta$.
Comme $Y$ est décent et que $X' \to Y$ est séparé, $X'$ est décent
(Espaces décents, lemmes
\ref{decent-spaces-lemma-properties-trivial-implications} et
\ref{decent-spaces-lemma-property-over-property}).
Ainsi, $f'$ est une altération d'après le
lemme \ref{lemma-alteration-generically-finite}.
\end{proof}








\section{Lieu schématique}
\label{section-schematic}

\noindent
Nous avons déjà démontré plusieurs résultats sur le lieu schématique
d'un espace algébrique. En voici une liste de références :
\begin{enumerate}
\item Propriétés des espaces, sections
\ref{spaces-properties-section-schematic} et
\ref{spaces-properties-section-getting-a-scheme},
\item Espaces décents, section \ref{decent-spaces-section-schematic},
\item Propriétés des espaces, lemme
\ref{spaces-properties-lemma-point-like-spaces}
$\Leftarrow$
Espaces décents, lemme \ref{decent-spaces-lemma-decent-point-like-spaces}
$\Leftarrow$
Espaces décents, lemme \ref{decent-spaces-lemma-when-field},
\item Limites d'espaces, section \ref{spaces-limits-section-affine}, et
\item Limites d'espaces, section \ref{spaces-limits-section-representable}.
\end{enumerate}
Dans certains cas, certains types de morphismes d'espaces algébriques
sont automatiquement représentables, par exemple les morphismes
séparés et localement quasi-finis (Morphismes d'espaces, lemme
\ref{spaces-morphisms-lemma-locally-quasi-finite-separated-representable})
et les monomorphismes plats (Compléments sur les morphismes d'espaces, lemme
\ref{spaces-more-morphisms-lemma-flat-case}).
Dans la section \ref{section-schematic-and-field-extension},
nous étudierons le comportement du lieu
schématique par extension du corps de base.

\begin{lemma}
\label{lemma-locally-finite-type-dim-zero}
Soit $S$ un schéma. Soit $X$ un espace algébrique sur $S$.
Supposons que $X$ satisfasse au moins l'une des conditions suivantes :
\begin{enumerate}
\item $X$ est quasi-séparé et $\dim(X) = 0$,
\item $X$ est localement de type fini sur un corps $k$ et $\dim(X) = 0$,
\item $X$ est noethérien et $\dim(X) = 0$, ou
\item en ajouter d'autres ici.
\end{enumerate}
Alors $X$ est un schéma séparé et tout ouvert quasi-compact de $X$ est affine.
\end{lemma}

\begin{proof}
Si nous démontrons que tout ouvert quasi-compact de $X$ est affine, alors
$X$ est un schéma séparé. Nous pouvons donc supposer $X$ quasi-compact
et cherchons à montrer que $X$ est affine.
Les cas (2) et (3) découlent immédiatement du cas (1), mais nous donnerons des
démonstrations séparées de (2) et (3), qui utilisent beaucoup moins de théorie.

\medskip\noindent
Démonstration de (3). Soit $U$ un schéma affine et soit $U \to X$ un
morphisme \'etale. Posons $R = U \times_X U$. Les deux morphismes de
projection $s, t : R \to U$ sont des morphismes étales de schémas. D'après la
définition de Propriétés des espaces \ref{spaces-properties-definition-dimension},
nous avons $\dim(U) = 0$ et $\dim(R) = 0$.
Comme $R$ est un schéma localement noethérien de dimension $0$,
$R$ est une réunion disjointe de spectres
d'anneaux locaux artiniens
(Propriétés, lemme \ref{properties-lemma-locally-Noetherian-dimension-0}).
Puisque nous avons supposé $X$ noethérien (donc quasi-séparé), nous
concluons que $R$ est quasi-compact. Ainsi, $R$ est un schéma affine
(utiliser Schémas, lemme \ref{schemes-lemma-disjoint-union-affines}).
Les morphismes étales $s, t : R \to U$ induisent des extensions finies
des corps résiduels. Le lemme d'Algèbre
\ref{algebra-lemma-essentially-of-finite-type-into-artinian-local}
montre donc que $s$ et $t$ sont finis
(nous omettons un détail mineur).
La proposition de Groupoïdes
\ref{groupoids-proposition-finite-flat-equivalence}
montre alors que $X = U/R$ est un schéma affine.

\medskip\noindent
Démonstration de (2), presque identique à celle de (3).
Soit $U$ un schéma affine et soit $U \to X$ un morphisme étale surjectif.
Posons $R = U \times_X U$. Les deux morphismes de projection
$s, t : R \to U$ sont des morphismes étales de schémas. D'après la
définition de Propriétés des espaces \ref{spaces-properties-definition-dimension},
nous avons $\dim(U) = 0$ et, de même, $\dim(R) = 0$.
D'autre part, le morphisme $U \to \Spec(k)$ est localement de type
fini comme composé du morphisme étale $U \to X$ et de
$X \to \Spec(k)$ ; voir
Morphismes d'espaces,
les lemmes \ref{spaces-morphisms-lemma-composition-finite-type} et
\ref{spaces-morphisms-lemma-etale-locally-finite-type}.
De même, $R \to \Spec(k)$ est localement de type fini.
Le lemme de Variétés
\ref{varieties-lemma-algebraic-scheme-dim-0}
montre donc que $U$ et $R$ sont des réunions disjointes de spectres de
$k$-algèbres locales artiniennes finies sur $k$. Il en va donc
de même pour $U \times_{\Spec(k)} U$. Comme
$$
R = U \times_X U \longrightarrow U \times_{\Spec(k)} U
$$
est un monomorphisme, $R$ est une réunion finie (!) de spectres de
$k$-algèbres finies. Il s'ensuit que $R$ est affine ; voir
Schémas, lemme \ref{schemes-lemma-disjoint-union-affines}.
En appliquant de nouveau le
lemme de Variétés \ref{varieties-lemma-algebraic-scheme-dim-0},
nous voyons que $R$ est fini sur $k$. Ainsi, $s, t$
sont finis ; voir
Morphismes, lemme \ref{morphisms-lemma-finite-permanence}.
La proposition de Groupoïdes
\ref{groupoids-proposition-finite-flat-equivalence}
montre alors que $X = U/R$ est un schéma affine.

\medskip\noindent
Démonstration cohomologique de (1). D'après le lemme de Cohomologie des espaces
\ref{spaces-cohomology-lemma-vanishing-above-dimension},
les groupes de cohomologie supérieure s'annulent pour tout
faisceau quasi-cohérent $\mathcal{F}$ sur $X$. Par conséquent, $X$
est affine (donc, en particulier, un schéma) d'après la
proposition de Cohomologie des espaces
\ref{spaces-cohomology-proposition-vanishing-affine}.

\medskip\noindent
Démonstration géométrique de (1). Choisissons une stratification
$$
\emptyset = U_{n + 1} \subset
U_n \subset U_{n - 1} \subset \ldots \subset U_1 = X
$$
et des morphismes étales $f_p : V_p \to U_p$ comme dans le
lemme d'Espaces décents
\ref{decent-spaces-lemma-filter-quasi-compact-quasi-separated}
(nous utiliserons ci-dessous toutes leurs propriétés).
Alors $\dim(V_p) = 0$ par notre définition de la dimension des espaces
algébriques. Le lemme de Propriétés \ref{properties-lemma-dimension-zero}
s'applique donc à chaque $V_p$. Puis $f_p^{-1}(U_{p + 1}) \subset V_p$
est un ouvert quasi-compact, donc à la fois affine et fermé.
Il s'ensuit que $|T_p| \subset |U_p|$ (voir la référence citée)
est à la fois ouvert et fermé. Ainsi, $X$ est une réunion disjointe
de sous-espaces ouverts et fermés dont les structures réduites sont des schémas.
Il s'ensuit que $X$ est un schéma
(Limites d'espaces, lemme \ref{spaces-limits-lemma-reduction-scheme}).
La démonstration s'achève alors par le cas des schémas
déjà cité ci-dessus.
\end{proof}

\noindent
Le lemme suivant affirme qu'un espace algébrique quasi-séparé
est un schéma hors d'une partie de codimension au moins $1$.

\begin{lemma}
\label{lemma-generic-point-in-schematic-locus}
Soit $S$ un schéma. Soit $X$ un espace algébrique quasi-séparé sur $S$.
Soit $x \in |X|$. Les conditions suivantes sont équivalentes :
\begin{enumerate}
\item $x$ est un point de codimension $0$ sur $X$ ;
\item l'anneau local de $X$ en $x$ est de dimension $0$ ;
\item $x$ est un point générique d'une composante irréductible de $|X|$.
\end{enumerate}
Si ces conditions sont satisfaites, il existe un sous-espace ouvert de $X$
qui contient $x$ et qui est un schéma.
\end{lemma}

\begin{proof}
L'équivalence de (1), (2) et (3) découle du
lemme d'Espaces décents \ref{decent-spaces-lemma-decent-generic-points}
et du fait qu'un espace algébrique quasi-séparé est décent
(Espaces décents, section \ref{decent-spaces-section-reasonable-decent}).
Nous donnerons toutefois, dans le paragraphe suivant, une démonstration plus
élémentaire de cette équivalence.

\medskip\noindent
Notons que (1) et (2) sont équivalents par définition
(Propriétés des espaces, définition
\ref{spaces-properties-definition-dimension-local-ring}).
Pour démontrer l'équivalence de (1) et (3), nous pouvons supposer $X$ quasi-compact.
Choisissons
$$
\emptyset = U_{n + 1} \subset
U_n \subset U_{n - 1} \subset \ldots \subset U_1 = X
$$
et $f_i : V_i \to U_i$ comme dans le lemme d'Espaces décents
\ref{decent-spaces-lemma-filter-quasi-compact-quasi-separated}.
Supposons $x \in U_i$, $x \not \in U_{i + 1}$. Alors $x = f_i(y)$ pour
un unique $y \in V_i$. Si (1) est satisfaite, $y$ est un point générique
d'une composante irréductible de $V_i$ (Propriétés des espaces, lemme
\ref{spaces-properties-lemma-codimension-0-points}).
Comme $f_i^{-1}(U_{i + 1})$ est un ouvert quasi-compact de $V_i$
qui ne contient pas $y$, il existe un voisinage ouvert $W \subset V_i$
de $y$ disjoint de $f_i^{-1}(V_i)$
(voir
Propriétés, lemme \ref{properties-lemma-generic-point-in-constructible},
ou, plus simplement, Algèbre, lemme
\ref{algebra-lemma-standard-open-containing-maximal-point}).
Alors $f_i|_W : W \to X$ est un isomorphisme sur son image, et donc
$x = f_i(y)$ est un point générique de $|X|$. Réciproquement, supposons (3).
Alors $f_i$ envoie $\overline{\{y\}}$ sur la composante irréductible
$\overline{\{x\}}$ de $|U_i|$. Comme $|f_i|$ est bijectif au-dessus de
$\overline{\{x\}}$, il s'ensuit que $\overline{\{y\}}$
est une composante irréductible de $U_i$. Ainsi, $x$ est un point de
codimension $0$.

\medskip\noindent
La dernière assertion du lemme est la
proposition de Propriétés des espaces
\ref{spaces-properties-proposition-locally-quasi-separated-open-dense-scheme}.
\end{proof}

\noindent
Le lemme suivant affirme qu'un espace algébrique séparé localement noethérien
est un schéma en codimension $1$, c'est-à-dire hors de la codimension $2$.

\begin{lemma}
\label{lemma-codim-1-point-in-schematic-locus}
\begin{slogan}
Les espaces algébriques séparés sont des schémas en codimension 1.
\end{slogan}
Soit $S$ un schéma. Soit $X$ un espace algébrique sur $S$.
Soit $x \in |X|$. Si $X$ est séparé et localement noethérien, et si
la dimension de l'anneau local de $X$ en $x$ est $\leq 1$
(Propriétés des espaces, définition
\ref{spaces-properties-definition-dimension-local-ring}),
alors il existe un sous-espace ouvert de $X$ contenant $x$ qui est un schéma.
\end{lemma}

\begin{proof}
(Voir la remarque ci-dessous pour une autre approche qui évite les résultats sur
les groupoïdes finis.) Nous pouvons remplacer $X$ par un voisinage quasi-compact
de $x$ ; nous pouvons donc supposer $X$ quasi-compact, séparé et noethérien.
Il existe un schéma $U$ et un morphisme fini surjectif $U \to X$ ;
voir Limites d'espaces, proposition
\ref{spaces-limits-proposition-there-is-a-scheme-finite-over}.
Soit $R = U \times_X U$. Alors $j : R \to U \times_S U$ définit une relation
d'équivalence et nous obtenons un schéma en groupoïdes $(U, R, s, t, c)$ sur $S$,
avec $s, t$ finis et $U$ noethérien et séparé.
Soit $\{u_1, \ldots, u_n\} \subset U$ l'ensemble des points ayant pour image $x$.
Alors $\dim(\mathcal{O}_{U, u_i}) \leq 1$ d'après le
lemme d'Espaces décents
\ref{decent-spaces-lemma-dimension-local-ring-quasi-finite}.

\medskip\noindent
D'après le lemme de Compléments sur les groupoïdes
\ref{more-groupoids-lemma-find-affine-codimension-1},
il existe un ouvert affine $R$-invariant $W \subset U$ qui contient
l'orbite $\{u_1, \ldots, u_n\}$. Comme $U \to X$ est fini surjectif,
l'application continue $|U| \to |X|$ est fermée et surjective, donc
submersive d'après le lemme de Topologie
\ref{topology-lemma-closed-morphism-quotient-topology}.
Ainsi, $f(W)$ est ouvert et il existe un sous-espace ouvert $X' \subset X$
tel que $f : W \to X'$ soit un morphisme fini surjectif.
Alors $X'$ est un schéma affine d'après le
lemme de Cohomologie des espaces
\ref{spaces-cohomology-lemma-image-affine-finite-morphism-affine-Noetherian},
ce qui achève la démonstration.
\end{proof}

\begin{remark}
\label{remark-alternate-proof-scheme-codim-1}
Voici une esquisse de démonstration du
lemme \ref{lemma-codim-1-point-in-schematic-locus}
qui n'utilise pas le
lemme de Compléments sur les groupoïdes
\ref{more-groupoids-lemma-find-affine-codimension-1}.

\medskip\noindent
Étape 1. Nous pouvons supposer que $X$ est un espace algébrique réduit,
noethérien et séparé (par exemple d'après Cohomologie des espaces, lemme
\ref{spaces-cohomology-lemma-image-affine-finite-morphism-affine-Noetherian},
ou d'après
Limites d'espaces, lemme \ref{spaces-limits-lemma-reduction-scheme}),
et nous pouvons choisir un morphisme fini surjectif
$Y \to X$, où $Y$ est un schéma noethérien (d'après
Limites d'espaces, proposition
\ref{spaces-limits-proposition-there-is-a-scheme-finite-over}).

\medskip\noindent
Étape 2. Après avoir remplacé $X$ par un voisinage ouvert de $x$, il
existe un morphisme fini birationnel $X' \to X$ et un sous-schéma fermé
$Y' \subset X' \times_X Y$ tels que $Y' \to X'$ soit surjectif,
fini et localement libre. En effet, puisque $X$ est réduit, il existe un
sous-espace ouvert dense $U \subset X$ au-dessus duquel $Y$ est plat (Morphismes d'espaces,
proposition \ref{spaces-morphisms-proposition-generic-flatness-reduced}).
Nous pouvons alors choisir un éclatement $U$-admissible $b : \tilde X \to X$ tel
que la transformée stricte $\tilde Y$ de $Y$ soit plate sur $\tilde X$ ; voir
Compléments sur les morphismes d'espaces, lemme
\ref{spaces-more-morphisms-lemma-flat-after-blowing-up}.
(On peut aussi utiliser les schémas de Hilbert si l'on veut éviter
le résultat sur les éclatements.)
Soit alors $X' \subset \tilde X$ l'adhérence
schématique de $b^{-1}(U)$ et $Y' = X' \times_{\tilde X} \tilde Y$.
Comme $x$ est un point de codimension $1$, $X' \to X$ est fini au-dessus d'un
voisinage de $x$ (lemme \ref{lemma-finite-in-codim-1}).

\medskip\noindent
Étape 3. Après avoir restreint $X$ à un voisinage plus petit de $x$, nous
obtenons que $X'$ est un schéma. En effet, $Y'$ est un schéma, $Y' \to X'$
est fini localement libre, et tout ensemble fini de points de codimension $1$
de $Y'$ est contenu dans un ouvert affine. Utiliser la
proposition de Propriétés des espaces
\ref{spaces-properties-proposition-finite-flat-equivalence-global}
et la
proposition de Variétés
\ref{varieties-proposition-finite-set-of-points-of-codim-1-in-affine}.

\medskip\noindent
Étape 4. Il existe un ouvert affine $W' \subset X'$ qui contient tous les points
situés au-dessus de $x$ et qui est l'image réciproque d'un sous-espace ouvert de $X$.
Pour le démontrer, soit $Z \subset X$ l'adhérence de l'ensemble des points
où $X' \to X$ n'est pas un isomorphisme. Nous pouvons supposer $x \in Z$, sinon
nous avons déjà terminé. Alors $x$ est un point générique d'une composante
irréductible de $Z$ et, après restriction de $X$, nous pouvons supposer $Z$ affine
(lemme \ref{lemma-generic-point-in-schematic-locus}).
L'image réciproque $Z' \subset X'$ est alors elle aussi un schéma affine.
Soient $x_1, \ldots, x_n \in Z'$ les points ayant pour image $x$.
Nous pouvons trouver un ouvert affine $W'$ de $X'$ dont l'intersection avec
$Z'$ est l'image réciproque d'un ouvert principal de $Z$ contenant $x$.
En effet, choisissons d'abord un ouvert affine $W' \subset X'$ contenant
$x_1, \ldots, x_n$ grâce à la proposition de Variétés
\ref{varieties-proposition-finite-set-of-points-of-codim-1-in-affine}.
Choisissons ensuite un ouvert principal $D(f) \subset Z$ contenant $x$
dont l'image réciproque $D(f|_{Z'})$ est contenue dans $W' \cap Z'$.
Choisissons enfin $f' \in \Gamma(W', \mathcal{O}_{W'})$ dont la restriction
est $f|_{Z'}$, puis remplaçons $W'$ par $D(f') \subset W'$.
Comme $X' \to X$ est un isomorphisme hors de $Z' \to Z$, le choix
de $W'$ garantit que l'image $W \subset X$ de $W'$ est ouverte
et qu'elle a pour image réciproque $W'$ dans $X'$.

\medskip\noindent
Étape 5. Alors $W' \to W$ est un morphisme fini surjectif et $W$ est un schéma d'après le
lemme de Cohomologie des espaces
\ref{spaces-cohomology-lemma-image-affine-finite-morphism-affine-Noetherian},
ce qui achève la démonstration.
\end{remark}





\section{Lieu schématique et extension de corps}
\label{section-schematic-and-field-extension}

\noindent
Il peut arriver qu'un espace algébrique non représentable sur un corps $k$
devienne représentable (c'est-à-dire un schéma) après changement de base à une extension
de $k$. Voir Espaces, exemple \ref{spaces-example-non-representable-descent}.
Dans cette section, nous étudions cette question.

\begin{lemma}
\label{lemma-scheme-after-purely-inseparable-base-change}
Soit $k$ un corps. Soit $X$ un espace algébrique sur $k$.
S'il existe une extension radicielle de corps $k'/k$
telle que $X_{k'}$ soit un schéma, alors $X$ est un schéma.
\end{lemma}

\begin{proof}
Le morphisme $X_{k'} \to X$ est entier, surjectif et
universellement injectif. Le lemme découle donc du
lemme de Limites d'espaces
\ref{spaces-limits-lemma-integral-universally-bijective-scheme}.
\end{proof}

\begin{lemma}
\label{lemma-when-scheme-after-base-change}
Soit $k$ un corps muni d'une clôture algébrique $\overline{k}$.
Soit $X$ un espace algébrique quasi-séparé sur $k$.
\begin{enumerate}
\item S'il existe une extension de corps $K/k$ telle que
$X_K$ soit un schéma, alors $X_{\overline{k}}$ est un schéma.
\item Si $X$ est quasi-compact et s'il existe une extension de corps
$K/k$ telle que $X_K$ soit un schéma, alors $X_{k'}$
est un schéma pour une certaine extension finie séparable $k'$ de $k$.
\end{enumerate}
\end{lemma}

\begin{proof}
Puisque tout espace algébrique est la réunion de ses sous-espaces ouverts
quasi-compacts, la première partie du lemme découle de
la seconde (nous omettons certains détails). Supposons donc $X$ quasi-compact
et donnée une extension $K/k$ pour laquelle $X_K$ est représentable.
Écrivons $K = \bigcup A$ comme la limite inductive des sous-$k$-algèbres de type fini
$A$. D'après le lemme de Limites d'espaces \ref{spaces-limits-lemma-limit-is-scheme},
$X_A$ est un schéma pour une certaine $A$. Choisissons un idéal maximal
$\mathfrak m \subset A$. D'après le théorème des zéros de Hilbert
(Algèbre, théorème \ref{algebra-theorem-nullstellensatz}),
le corps résiduel $k' = A/\mathfrak m$ est une extension finie de $k$.
Ainsi, $X_{k'}$ est un schéma. Si $k' \supset k$ n'est pas
séparable, soit $k'/k''/k$ la sous-extension
fournie par le lemme de Corps \ref{fields-lemma-separable-first}.
Comme $k'/k''$ est radicielle, le
lemme \ref{lemma-scheme-after-purely-inseparable-base-change}
montre que l'espace algébrique $X_{k''}$ est un schéma. Puisque $k''|k$ est séparable,
la démonstration est achevée.
\end{proof}

\begin{lemma}
\label{lemma-base-change-by-Galois}
Soit $k'/k$ une extension galoisienne finie de groupe de Galois $G$.
Soit $X$ un espace algébrique sur $k$. Alors $G$ agit librement sur
l'espace algébrique $X_{k'}$ et $X = X_{k'}/G$ au sens du
lemme de Propriétés des espaces \ref{spaces-properties-lemma-quotient}.
\end{lemma}

\begin{proof}
Démonstration omise. Indications : montrer d'abord que $\Spec(k) = \Spec(k')/G$.
Utiliser ensuite la compatibilité des quotients avec le changement de base.
\end{proof}

\begin{lemma}
\label{lemma-when-quotient-scheme-at-point}
Soit $S$ un schéma. Soit $X$ un espace algébrique sur $S$ et
soit $G$ un groupe fini agissant librement sur $X$. Posons $Y = X/G$ comme
dans le lemme de Propriétés des espaces \ref{spaces-properties-lemma-quotient}.
Pour $y \in |Y|$, les conditions suivantes sont équivalentes :
\begin{enumerate}
\item $y$ appartient au lieu schématique de $Y$ ;
\item il existe un ouvert affine $U \subset X$
qui contient l'image réciproque de $y$.
\end{enumerate}
\end{lemma}

\begin{proof}
Il découle de la construction de $Y = X/G$ dans le
lemme de Propriétés des espaces \ref{spaces-properties-lemma-quotient}
que le morphisme $X \to Y$ est surjectif et \'etale.
Nous avons bien sûr $X \times_Y X = X \times G$ ; le morphisme
$X \to Y$ est donc même fini étale. Il est aussi surjectif.
Le lemme découle ainsi du
lemme d'Espaces décents \ref{decent-spaces-lemma-when-quotient-scheme-at-point}.
\end{proof}

\begin{lemma}
\label{lemma-scheme-after-purely-transcendental-base-change}
Soit $k$ un corps. Soit $X$ un espace algébrique
quasi-séparé sur $k$. S'il existe une extension transcendante pure
$K/k$ telle que $X_K$ soit un schéma, alors
$X$ est un schéma.
\end{lemma}

\begin{proof}
Puisque tout espace algébrique est la réunion de ses sous-espaces ouverts
quasi-compacts, nous pouvons supposer $X$ quasi-compact (certains détails sont omis).
Rappelons (Corps, définition \ref{fields-definition-transcendence})
que l'hypothèse sur l'extension $K/k$ signifie que
$K$ est le corps des fractions d'un anneau de polynômes (en un nombre éventuellement infini
de variables) sur $k$. Ainsi, $K = \bigcup A$ est la réunion de sous-algèbres
dont chacune est une localisation d'une algèbre polynomiale en un nombre fini de variables sur $k$.
D'après le lemme de Limites d'espaces \ref{spaces-limits-lemma-limit-is-scheme},
$X_A$ est un schéma pour une certaine $A$. Écrivons
$$
A = k[x_1, \ldots, x_n][1/f]
$$
pour un certain élément non nul $f \in k[x_1, \ldots, x_n]$.

\medskip\noindent
Si $k$ est infini, nous pouvons achever la démonstration comme suit : choisissons
$a_1, \ldots, a_n \in k$ tels que $f(a_1, \ldots, a_n) \not = 0$.
Alors $(a_1, \ldots, a_n)$ définit un homomorphisme de $k$-algèbres $A \to k$
qui envoie $x_i$ sur $a_i$ et $1/f$ sur $1/f(a_1, \ldots, a_n)$.
Ainsi, le changement de base $X_A \times_{\Spec(A)} \Spec(k) \cong X$ est un
schéma, comme voulu.

\medskip\noindent
Dans ce paragraphe, achevons la démonstration lorsque $k$ est fini. Dans ce
cas, écrivons $X = \lim X_i$, où les $X_i$ sont de présentation finie sur $k$
et les morphismes de transition sont affines
(Limites d'espaces, lemme \ref{spaces-limits-lemma-relative-approximation}).
En utilisant le lemme de Limites d'espaces \ref{spaces-limits-lemma-limit-is-scheme},
nous voyons que $X_{i, A}$ est un schéma pour un certain $i$. Nous pouvons donc supposer
$X \to \Spec(k)$ de présentation finie. Soit $x \in |X|$ un point fermé.
Nous pouvons représenter $x$ par une immersion fermée
$\Spec(\kappa) \to X$
(Espaces décents, lemme \ref{decent-spaces-lemma-decent-space-closed-point}).
Alors $\Spec(\kappa) \to \Spec(k)$ est de type fini, donc $\kappa$
est une extension finie de $k$ (d'après le théorème des zéros de Hilbert ; voir
Algèbre, théorème \ref{algebra-theorem-nullstellensatz} ;
nous omettons certains détails). Posons $[\kappa : k] = d$. Choisissons un entier
$n \gg 0$ premier à $d$ et soit $k'/k$ l'extension
de degré $n$. Alors $k'/k$ est galoisienne, avec $G = \text{Aut}(k'/k)$
cyclique d'ordre $n$. Si $n$ est assez grand, il existe un homomorphisme de $k$-algèbres
$A \to k'$ pour la même raison que ci-dessus.
Alors $X_{k'}$ est un schéma et $X = X_{k'}/G$
(lemme \ref{lemma-base-change-by-Galois}).
D'autre part, comme $n$ et $d$ sont premiers entre eux, nous avons
$$
\Spec(\kappa) \times_{X} X_{k'} =
\Spec(\kappa) \times_{\Spec(k)} \Spec(k') =
\Spec(\kappa \otimes_k k')
$$
est le spectre d'un corps. Autrement dit, la fibre de $X_{k'} \to X$
au-dessus de $x$ est réduite à un seul point. Le
lemme \ref{lemma-when-quotient-scheme-at-point}
montre donc que $x$ appartient au lieu schématique de $X$, comme voulu.
\end{proof}

\begin{remark}
\label{remark-when-does-the-argument-work}
Soit $k$ un corps fini. Soit $K/k$ une extension de corps géométriquement
irréductible. Alors $K$ est la limite inductive de $k$-algèbres de type fini
géométriquement irréductibles $A$. Pour une telle $A$, les estimations
de Lang et Weil \cite{LW} montrent que, pour $n \gg 0$, il existe
un homomorphisme de $k$-algèbres $A \to k'$ avec $k'/k$ de degré $n$.
L'analyse de l'argument donné dans la démonstration du
lemme \ref{lemma-scheme-after-purely-transcendental-base-change}
montre que, si $X$ est un espace algébrique quasi-séparé sur $k$
et si $X_K$ est un schéma, alors $X$ est un schéma. Si nous avons un jour besoin de ce
résultat, nous le formulerons précisément et le démontrerons ici.
\end{remark}

\begin{lemma}
\label{lemma-scheme-over-algebraic-closure-enough-affines}
Soit $k$ un corps muni d'une clôture algébrique $\overline{k}$. Soit $X$
un espace algébrique sur $k$ tel que
\begin{enumerate}
\item $X$ est décent et localement de type fini sur $k$ ;
\item $X_{\overline{k}}$ est un schéma ;
\item tout ensemble fini de points $\overline{k}$-rationnels de $X_{\overline{k}}$
est contenu dans un ouvert affine.
\end{enumerate}
Alors $X$ est un schéma.
\end{lemma}

\begin{proof}
Si $K/k$ est une extension, le changement de base $X_K$ est
décent (Espaces décents, lemme
\ref{decent-spaces-lemma-representable-named-properties})
et localement de type fini
sur $K$ (Morphismes d'espaces, lemme
\ref{spaces-morphisms-lemma-base-change-finite-type}).
D'après le lemme \ref{lemma-scheme-after-purely-inseparable-base-change},
il suffit de démontrer que $X$ devient un schéma après changement de base à
la perfection de $k$ ; nous pouvons donc supposer que $k$ est un corps parfait
(cette étape n'est pas strictement nécessaire, mais facilite
la compréhension des autres arguments).
En recouvrant $X$ par des ouverts quasi-compacts, il suffit de démontrer
le lemme lorsque $X$ est quasi-compact (nous omettons un détail mineur).
Dans ce cas, $|X|$ est un espace topologique sobre
(Espaces décents, proposition
\ref{decent-spaces-proposition-reasonable-sober}).
Il suffit donc de montrer que tout point fermé de $|X|$
appartient au lieu schématique de $X$
(utiliser Propriétés des espaces, lemme \ref{spaces-properties-lemma-subscheme}, et
Topologie, lemme \ref{topology-lemma-quasi-compact-closed-point}).

\medskip\noindent
Soit $x \in |X|$ un point fermé. D'après le lemme d'Espaces décents
\ref{decent-spaces-lemma-decent-space-closed-point},
il existe une immersion fermée $\Spec(l) \to X$ représentant $x$.
Alors $\Spec(l) \to \Spec(k)$ est de type fini (Morphismes d'espaces,
lemme \ref{spaces-morphisms-lemma-composition-finite-type}) et nous
concluons que $l$ est une extension finie de $k$
d'après le théorème des zéros de Hilbert (Algèbre, théorème
\ref{algebra-theorem-nullstellensatz}). Elle est séparable, car
$k$ est parfait. Ainsi, le schéma
$$
\Spec(l) \times_X X_{\overline{k}} =
\Spec(l) \times_{\Spec(k)} \Spec(\overline{k}) =
\Spec(l \otimes_k \overline{k})
$$
est la réunion disjointe d'un nombre fini de points $\overline{k}$-rationnels.
D'après l'hypothèse (3), il existe un ouvert affine $W \subset X_{\overline{k}}$
qui contient ces points.

\medskip\noindent
D'après le lemme \ref{lemma-when-scheme-after-base-change}, $X_{k'}$
est un schéma pour une certaine extension finie $k'/k$. Après avoir agrandi
$k'$, nous pouvons supposer qu'il existe un ouvert affine $U' \subset X_{k'}$
dont le changement de base à $\overline{k}$ redonne $W$
(utiliser que $X_{\overline{k}}$ est la limite des schémas $X_{k''}$
pour les extensions finies $k' \subset k'' \subset \overline{k}$, ainsi que les
lemmes de Limites \ref{limits-lemma-descend-opens} et
\ref{limits-lemma-limit-affine}). Nous pouvons supposer
que $k'/k$ est une extension galoisienne (prendre la clôture normale ;
Corps, lemme \ref{fields-lemma-normal-closure}, et utiliser
que $k$ est parfait). Posons $G = \text{Gal}(k'/k)$.
Par construction, le sous-schéma fermé stable par $G$
$\Spec(l) \times_X X_{k'}$ est contenu dans $U'$.
Ainsi, $x$ appartient au lieu schématique d'après les
lemmes \ref{lemma-base-change-by-Galois} et
\ref{lemma-when-quotient-scheme-at-point}.
\end{proof}

\noindent
Les deux lemmes suivants devraient figurer ailleurs.
Comparer le lemme suivant au
lemme d'Espaces décents \ref{decent-spaces-lemma-conditions-on-space-over-field}.

\begin{lemma}
\label{lemma-locally-quasi-finite-over-field}
Soit $k$ un corps. Soit $X$ un espace algébrique sur $k$.
Les conditions suivantes sont équivalentes :
\begin{enumerate}
\item $X$ est localement quasi-fini sur $k$ ;
\item $X$ est localement de type fini sur $k$ et de dimension $0$ ;
\item $X$ est un schéma localement quasi-fini sur $k$ ;
\item $X$ est un schéma localement de type fini sur $k$ et de
dimension $0$ ;
\item $X$ est une réunion disjointe de spectres de $k$-algèbres locales artiniennes
$A$ finies sur $k$ telles que $\dim_k(A) < \infty$.
\end{enumerate}
\end{lemma}

\begin{proof}
Comme nous sommes sur un corps, la dimension relative de $X/k$ est égale à
la dimension de $X$. Le
lemme de Morphismes d'espaces
\ref{spaces-morphisms-lemma-locally-quasi-finite-rel-dimension-0}
montre donc que (1) et (2) sont équivalents. Il découle alors du
lemme \ref{lemma-locally-finite-type-dim-zero}
(et d'implications immédiates) que (1) -- (4) sont équivalents.
Enfin, le
lemme de Variétés \ref{varieties-lemma-algebraic-scheme-dim-0}
montre que (1) -- (4) sont équivalents à (5).
\end{proof}

\begin{lemma}
\label{lemma-mono-towards-locally-quasi-finite-over-field}
Soit $k$ un corps. Soit $f : X \to Y$ un monomorphisme d'espaces algébriques
sur $k$. Si $Y$ est localement quasi-fini sur $k$, alors $X$ l'est aussi.
\end{lemma}

\begin{proof}
Supposons $Y$ localement quasi-fini sur $k$. D'après le
lemme \ref{lemma-locally-quasi-finite-over-field},
nous avons $Y = \coprod \Spec(A_i)$, où chaque $A_i$ est un
anneau local artinien fini sur $k$. D'après le
lemme d'Espaces décents
\ref{decent-spaces-lemma-monomorphism-toward-disjoint-union-dim-0-rings},
$X$ est un schéma. Considérons $X_i = f^{-1}(\Spec(A_i))$.
Alors $X_i$ possède soit un point, soit aucun. Si $X_i$ n'a aucun point, il
n'y a rien à démontrer. Si $X_i$ possède un point, alors
$X_i = \Spec(B_i)$, où $B_i$ est un anneau local de dimension zéro,
et $A_i \to B_i$ est un épimorphisme d'anneaux. En particulier,
$A_i/\mathfrak m_{A_i} = B_i/\mathfrak m_{A_i}B_i$, et
$A_i \to B_i$ est surjectif d'après le lemme de Nakayama,
Algèbre, lemme \ref{algebra-lemma-NAK}
(car $\mathfrak m_{A_i}$ est un idéal nilpotent !).
Ainsi, $B_i$ est une $k$-algèbre locale finie, et le
lemme \ref{lemma-locally-quasi-finite-over-field}
montre que $X \to \Spec(k)$ est localement quasi-fini.
\end{proof}











\section{Espaces algébriques géométriquement réduits}
\label{section-geometrically-reduced}

\noindent
Si $X$ est un espace algébrique réduit sur un corps, il peut arriver que $X$
cesse d'être réduit après extension du corps de base. Cela ne se produit pas
pour les espaces algébriques géométriquement réduits.

\begin{definition}
\label{definition-geometrically-reduced}
Soit $k$ un corps.
Soit $X$ un espace algébrique sur $k$.
\begin{enumerate}
\item Soit $x \in |X|$ un point. Nous disons que $X$ est
{\it géométriquement réduit en $x$} si $\mathcal{O}_{X, \overline{x}}$
est géométriquement réduit sur $k$.
\item Nous disons que $X$ est {\it géométriquement réduit} sur $k$
si $X$ est géométriquement réduit en tout point de $X$.
\end{enumerate}
\end{definition}

\noindent
Observons que, si $X$ est géométriquement réduit en $x$, alors l'anneau local
de $X$ en $x$ est réduit (Propriétés des espaces, lemme
\ref{spaces-properties-lemma-reduced-local-ring}).
De même, si $X$ est géométriquement réduit sur $k$, alors $X$ est réduit
(d'après Propriétés des espaces, lemme \ref{spaces-properties-lemma-reduced-space}).
Le lemme suivant implique notamment que cette définition est compatible
avec la propriété correspondante pour les schémas sur un corps.

\begin{lemma}
\label{lemma-geometrically-reduced-at-point}
Soit $k$ un corps. Soit $X$ un espace algébrique sur $k$.
Soit $x \in |X|$. Les conditions suivantes sont équivalentes :
\begin{enumerate}
\item $X$ est géométriquement réduit en $x$ ;
\item il existe un voisinage \'etale $(U, u) \to (X, x)$
tel que $U$ soit un schéma et que $U$ soit géométriquement réduit en $u$ ;
\item pour tout voisinage \'etale $(U, u) \to (X, x)$
tel que $U$ soit un schéma, $U$ est géométriquement réduit en $u$.
\end{enumerate}
\end{lemma}

\begin{proof}
Rappelons que l'anneau local $\mathcal{O}_{X, \overline{x}}$
est l'hensélisé strict de $\mathcal{O}_{U, u}$ ; voir le
lemme de Propriétés des espaces
\ref{spaces-properties-lemma-describe-etale-local-ring}.
D'après le lemme de Variétés \ref{varieties-lemma-geometrically-reduced-at-point},
$U$ est géométriquement réduit en $u$ si et seulement
si $\mathcal{O}_{U, u}$ est géométriquement réduit sur $k$.
Il nous reste donc à montrer ceci : si $A$ est une $k$-algèbre locale,
$A$ est géométriquement réduite sur $k$ si et seulement si
$A^{sh}$ est géométriquement réduite sur $k$.
Vérifions-le à partir de la définition des algèbres
géométriquement réduites (Algèbre, définition \ref{algebra-definition-geometrically-reduced}).
Soit $K/k$ une extension de corps.
Comme $A \to A^{sh}$ est fidèlement plat
(Compléments d'algèbre, lemme \ref{more-algebra-lemma-dumb-properties-henselization}),
le morphisme
$A \otimes_k K \to A^{sh} \otimes_k K$ est fidèlement plat
(Algèbre, lemme \ref{algebra-lemma-flat-base-change}).
Par conséquent, si $A^{sh} \otimes_k K$ est réduit, alors
$A \otimes_k K$ l'est aussi, d'après le lemme d'Algèbre \ref{algebra-lemma-descent-reduced}.
Réciproquement, rappelons que $A^{sh}$ est une limite inductive
d'$A$-algèbres \'etales ; voir le
lemme d'Algèbre \ref{algebra-lemma-strict-henselization}.
Ainsi, $A^{sh} \otimes_k K$ est une limite inductive
filtrante d'algèbres \'etales sur $A \otimes_k K$.
Nous concluons par le lemme d'Algèbre \ref{algebra-lemma-reduced-goes-up}.
\end{proof}

\begin{lemma}
\label{lemma-geometrically-reduced}
Soit $k$ un corps. Soit $X$ un espace algébrique sur $k$.
Les conditions suivantes sont équivalentes :
\begin{enumerate}
\item $X$ est géométriquement réduit ;
\item il existe un morphisme \'etale surjectif $U \to X$, où $U$
est un schéma, et $U$ est géométriquement réduit ;
\item pour tout morphisme \'etale $U \to X$
où $U$ est un schéma, $U$ est géométriquement réduit.
\end{enumerate}
\end{lemma}

\begin{proof}
Cela résulte immédiatement des définitions et du
lemme \ref{lemma-geometrically-reduced-at-point}.
\end{proof}

\noindent
Cette notion ne présente pas d'intérêt en caractéristique nulle.

\begin{lemma}
\label{lemma-perfect-reduced}
Soit $X$ un espace algébrique sur un corps parfait $k$ (par exemple,
$k$ est de caractéristique nulle).
\begin{enumerate}
\item Pour $x \in |X|$, si $\mathcal{O}_{X, \overline{x}}$ est
réduit, alors $X$ est géométriquement réduit en $x$.
\item Si $X$ est réduit, alors $X$ est géométriquement réduit sur $k$.
\end{enumerate}
\end{lemma}

\begin{proof}
La première assertion résulte du
lemme d'Algèbre \ref{algebra-lemma-separable-extension-preserves-reducedness}
et de la définition d'un corps parfait
(Algèbre, définition \ref{algebra-definition-perfect}).
La seconde assertion résulte de la première.
\end{proof}

\begin{lemma}
\label{lemma-geometrically-reduced-positive-characteristic}
Soit $k$ un corps de caractéristique $p > 0$. Soit $X$ un espace algébrique
sur $k$. Les conditions suivantes sont équivalentes :
\begin{enumerate}
\item $X$ est géométriquement réduit sur $k$ ;
\item $X_{k'}$ est réduit pour toute extension de corps $k'/k$ ;
\item $X_{k'}$ est réduit pour toute
extension radicielle finie de corps $k'/k$ ;
\item $X_{k^{1/p}}$ est réduit ;
\item $X_{k^{perf}}$ est réduit ;
\item $X_{\bar k}$ est réduit.
\end{enumerate}
\end{lemma}

\begin{proof}
Choisissons un morphisme \'etale surjectif $U \to X$, où $U$ est un schéma.
Le lemme \ref{lemma-geometrically-reduced} ramène l'énoncé au
résultat pour $U$ sur $k$.
Voir le lemme de Variétés \ref{varieties-lemma-geometrically-reduced}.
\end{proof}

\begin{lemma}
\label{lemma-geometrically-reduced-upstairs}
Soit $k$ un corps. Soit $X$ un espace algébrique sur $k$.
Soit $k'/k$ une extension de corps. Soit $x \in |X|$ un point et soit
$x' \in |X_{k'}|$ un point au-dessus de $x$.
Les conditions suivantes sont équivalentes :
\begin{enumerate}
\item $X$ est géométriquement réduit en $x$ ;
\item $X_{k'}$ est géométriquement réduit en $x'$.
\end{enumerate}
En particulier, $X$ est géométriquement réduit sur $k$ si et seulement si
$X_{k'}$ est géométriquement réduit sur $k'$.
\end{lemma}

\begin{proof}
Choisissons un morphisme \'etale $U \to X$, où $U$ est un schéma,
et un point $u \in U$ dont l'image est $x \in |X|$.
D'après le lemme de Propriétés des espaces \ref{spaces-properties-lemma-points-cartesian},
nous pouvons choisir un point $u' \in U_{k'} = U \times_X X_{k'}$ dont les images
sont $u$ et $x'$. D'après le lemme \ref{lemma-geometrically-reduced-at-point},
l'énoncé se ramène au lemme pour $U, u, u'$, à savoir au
lemme de Variétés \ref{varieties-lemma-geometrically-reduced-upstairs}.
\end{proof}

\begin{lemma}
\label{lemma-geometrically-reduced-etale-local}
Soit $k$ un corps. Soit $f : X \to Y$ un morphisme
d'espaces algébriques sur $k$. Soit $x \in |X|$ un point
dont l'image est $y \in |Y|$.
\begin{enumerate}
\item si $f$ est \'etale en $x$, alors
$X$ est géométriquement réduit en $x$ $\Leftrightarrow$
$Y$ est géométriquement réduit en $y$ ;
\item si $f$ est \'etale surjectif, alors
$X$ est géométriquement réduit $\Leftrightarrow$
$Y$ est géométriquement réduit.
\end{enumerate}
\end{lemma}

\begin{proof}
L'assertion (1) est claire, car
$\mathcal{O}_{X, \overline{x}} = \mathcal{O}_{Y, \overline{y}}$
lorsque $f$ est \'etale en $x$.
L'assertion (2) découle immédiatement de l'assertion (1).
\end{proof}






\section{Espaces algébriques géométriquement connexes}
\label{section-geometrically-connected}

\noindent
Si $X$ est un espace algébrique connexe sur un corps, il peut arriver que
$X$ devienne non connexe après extension du corps de base. Cela ne se
produit pas pour les espaces algébriques géométriquement connexes.

\begin{definition}
\label{definition-geometrically-connected}
Soit $X$ un espace algébrique sur le corps $k$. Nous disons que $X$ est
{\it géométriquement connexe} sur $k$ si le changement de base $X_{k'}$
est connexe pour toute extension de corps $k'$ de $k$.
\end{definition}

\noindent
Par convention, un espace topologique connexe est non vide ; à plus forte raison,
les espaces algébriques géométriquement connexes sont non vides.

\begin{lemma}
\label{lemma-geometrically-connected-check-after-extension}
Soit $X$ un espace algébrique sur le corps $k$.
Soit $k'/k$ une extension de corps.
Alors $X$ est géométriquement connexe sur $k$ si et seulement si
$X_{k'}$ est géométriquement connexe sur $k'$.
\end{lemma}

\begin{proof}
Si $X$ est géométriquement connexe sur $k$, il est clair que
$X_{k'}$ est géométriquement connexe sur $k'$. Réciproquement,
pour toute extension de corps $k''/k$, il existe une extension de corps commune
$k'''/k$ de $k'/k$ et $k''/k$ ; voir le
lemme de Corps \ref{fields-lemma-common-extension-field}. Comme le
morphisme $X_{k'''} \to X_{k''}$ est surjectif (c'est un changement de base
d'un morphisme surjectif entre spectres de corps), la
connexité de $X_{k'''}$ entraîne celle de $X_{k''}$.
Ainsi, si $X_{k'}$ est géométriquement connexe sur $k'$, alors
$X$ est géométriquement connexe sur $k$.
\end{proof}

\begin{lemma}
\label{lemma-bijection-connected-components}
Soit $k$ un corps. Soient $X$, $Y$ des espaces algébriques sur $k$.
Supposons $X$ géométriquement connexe sur $k$.
Alors le morphisme de projection
$$
p : X \times_k Y \longrightarrow Y
$$
induit une bijection entre les composantes connexes.
\end{lemma}

\begin{proof}
Soit $y \in |Y|$ représenté par un morphisme $\Spec(K) \to Y$,
où $K$ est un corps. La fibre de $|X \times_k Y| \to |Y|$ au-dessus de $y$
est l'image de $|X_K| \to |X \times_k Y|$ d'après le
lemme de Propriétés des espaces \ref{spaces-properties-lemma-points-cartesian}.
Ces fibres sont donc connexes puisque $X$ est, par hypothèse,
géométriquement connexe. D'après le
lemme de Morphismes d'espaces
\ref{spaces-morphisms-lemma-space-over-field-universally-open}
l'application $|p|$ est ouverte.
Nous pouvons donc appliquer le lemme de Topologie
\ref{topology-lemma-connected-fibres-connected-components}
pour conclure.
\end{proof}

\begin{lemma}
\label{lemma-separably-closed-field-connected-components}
Soit $k'/k$ une extension de corps. Soit $X$ un espace algébrique
sur $k$. Supposons $k$ séparablement clos. Alors le morphisme
$X_{k'} \to X$ induit une bijection entre les composantes connexes. En particulier,
$X$ est géométriquement connexe sur $k$ si et seulement si $X$ est connexe.
\end{lemma}

\begin{proof}
Comme $k$ est séparablement clos, nous voyons que
$k'$ est géométriquement connexe sur $k$ ; voir le
lemme d'Algèbre
\ref{algebra-lemma-separably-closed-connected-implies-geometric}.
Par conséquent, $Z = \Spec(k')$ est géométriquement connexe sur $k$ d'après le
lemme de Variétés \ref{varieties-lemma-affine-geometrically-connected}.
Comme $X_{k'} = Z \times_k X$, le résultat est un cas particulier du
lemme \ref{lemma-bijection-connected-components}.
\end{proof}

\begin{lemma}
\label{lemma-characterize-geometrically-connected}
Soit $k$ un corps. Soit $X$ un espace algébrique sur $k$.
Soit $\overline{k}$ une clôture séparable de $k$.
Alors $X$ est géométriquement connexe si et seulement si le changement de base
$X_{\overline{k}}$ est connexe.
\end{lemma}

\begin{proof}
Supposons $X_{\overline{k}}$ connexe. Soit $k'/k$ une extension de
corps. Il existe une extension de corps $\overline{k}'/\overline{k}$
dans laquelle $k'$ se plonge dans $\overline{k}'$ par-dessus $k$.
D'après le lemme \ref{lemma-separably-closed-field-connected-components},
$X_{\overline{k}'}$ est connexe.
Comme $X_{\overline{k}'} \to X_{k'}$ est surjectif, nous concluons
que $X_{k'}$ est connexe, comme voulu.
\end{proof}

\noindent
Soit $k$ un corps. Soit $\overline{k}/k$ une extension galoisienne
(éventuellement infinie). Par exemple, $\overline{k}$ pourrait être la
clôture séparable de $k$.
À tout $\sigma \in \text{Gal}(\overline{k}/k)$ correspond un
automorphisme
$
\Spec(\sigma) :
\Spec(\overline{k})
\longrightarrow
\Spec(\overline{k})
$.
Remarquons que
$\Spec(\sigma) \circ \Spec(\tau) = \Spec(\tau \circ \sigma)$.
Nous obtenons donc une action
$$
\text{Gal}(\overline{k}/k)^{opp} \times \Spec(\overline{k})
\longrightarrow
\Spec(\overline{k})
$$
du groupe opposé sur le schéma $\Spec(\overline{k})$.
Soit $X$ un espace algébrique sur $k$. Comme
$X_{\overline{k}} =
\Spec(\overline{k}) \times_{\Spec(k)} X$
par définition, l'action ci-dessus induit une action canonique
\begin{equation}
\label{equation-galois-action-base-change-kbar}
\text{Gal}(\overline{k}/k)^{opp} \times X_{\overline{k}}
\longrightarrow
X_{\overline{k}}.
\end{equation}

\begin{lemma}
\label{lemma-Galois-action-quasi-compact-open}
Soit $k$ un corps. Soit $X$ un espace algébrique sur $k$.
Soit $\overline{k}$ une extension galoisienne (éventuellement infinie) de $k$.
Soit $V \subset X_{\overline{k}}$ un ouvert quasi-compact.
Alors :
\begin{enumerate}
\item il existe une sous-extension finie $\overline{k}/k'/k$
et un ouvert quasi-compact $V' \subset X_{k'}$ tels que
$V = (V')_{\overline{k}}$,
\item il existe un sous-groupe ouvert $H \subset \text{Gal}(\overline{k}/k)$
tel que $\sigma(V) = V$ pour tout $\sigma \in H$.
\end{enumerate}
\end{lemma}

\begin{proof}
Choisissons un schéma $U$ et un morphisme \'etale surjectif $U \to X$.
Choisissons un ouvert quasi-compact $W \subset U_{\overline{k}}$ dont
l'image dans $X_{\overline{k}}$ est $V$. Cela est possible, car
$|U_{\overline{k}}| \to |X_{\overline{k}}|$ est continue et
$|U_{\overline{k}}|$ possède une base d'ouverts quasi-compacts. Nous pouvons appliquer le
lemme de Variétés
\ref{varieties-lemma-Galois-action-quasi-compact-open}
à $W \subset U_{\overline{k}}$, ce qui donne le lemme.
\end{proof}

\begin{lemma}
\label{lemma-closed-fixed-by-Galois}
Soit $k$ un corps. Soit $\overline{k}/k$ une extension galoisienne
(éventuellement infinie). Soit $X$ un espace algébrique sur $k$. Soit
$\overline{T} \subset |X_{\overline{k}}|$ vérifiant les propriétés suivantes :
\begin{enumerate}
\item $\overline{T}$ est une partie fermée de $|X_{\overline{k}}|$ ;
\item pour tout $\sigma \in \text{Gal}(\overline{k}/k)$,
nous avons $\sigma(\overline{T}) = \overline{T}$.
\end{enumerate}
Alors il existe une partie fermée $T \subset |X|$ dont l'image réciproque
dans $|X_{k'}|$ est $\overline{T}$.
\end{lemma}

\begin{proof}
Soit $T \subset |X|$ l'image de $\overline{T}$.
Comme $|X_{\overline{k}}| \to |X|$ est surjectif, l'assertion signifie
que $T$ est fermé et que son image réciproque est $\overline{T}$.
Choisissons un schéma $U$ et un morphisme \'etale surjectif $U \to X$.
Dans le cas des schémas
(voir le lemme de Variétés \ref{varieties-lemma-closed-fixed-by-Galois}),
il existe une partie fermée $T' \subset |U|$ dont l'image réciproque
dans $|U_{\overline{k}}|$ est l'image réciproque de $\overline{T}$.
Comme $|U_{\overline{k}}| \to |X_{\overline{k}}|$ est surjectif,
$T'$ est l'image réciproque de $T$ par $|U| \to |X|$.
Par notre construction de la topologie sur $|X|$, cela signifie que $T$ est
fermé. De la même manière, on voit que $\overline{T}$ est l'image
réciproque de $T$.
\end{proof}

\begin{lemma}
\label{lemma-characterize-geometrically-disconnected}
Soit $k$ un corps. Soit $X$ un espace algébrique sur $k$.
Les conditions suivantes sont équivalentes :
\begin{enumerate}
\item $X$ est géométriquement connexe ;
\item pour toute extension de corps finie séparable $k'/k$,
l'espace algébrique $X_{k'}$ est connexe.
\end{enumerate}
\end{lemma}

\begin{proof}
Cette démonstration est identique à celle du
lemme de Variétés \ref{varieties-lemma-characterize-geometrically-disconnected},
si ce n'est que
nous remplaçons le
lemme de Variétés \ref{varieties-lemma-characterize-geometrically-connected}
par le lemme \ref{lemma-characterize-geometrically-connected},
nous remplaçons le
lemme de Variétés \ref{varieties-lemma-Galois-action-quasi-compact-open}
par le lemme \ref{lemma-Galois-action-quasi-compact-open}, et
nous remplaçons le
lemme de Variétés \ref{varieties-lemma-closed-fixed-by-Galois}
par le lemme \ref{lemma-closed-fixed-by-Galois}.
Nous invitons le lecteur à lire cette démonstration plutôt que celle-ci.

\medskip\noindent
Il découle immédiatement de la définition que (1) implique (2).
Supposons que $X$ ne soit pas géométriquement connexe.
Soit $k \subset \overline{k}$ une clôture séparable
de $k$. D'après le
lemme \ref{lemma-characterize-geometrically-connected},
$X_{\overline{k}}$ est non connexe.
Écrivons $X_{\overline{k}} = \overline{U} \amalg \overline{V}$,
où $\overline{U}$ et $\overline{V}$ sont des sous-espaces algébriques ouverts, fermés et non vides
de $X_{\overline{k}}$.

\medskip\noindent
Soit $W \subset X$ un sous-espace ouvert quasi-compact quelconque.
Alors $W_{\overline{k}} \cap \overline{U}$ et
$W_{\overline{k}} \cap \overline{V}$ sont des sous-espaces ouverts et fermés de
$W_{\overline{k}}$. En particulier, $W_{\overline{k}} \cap \overline{U}$ et
$W_{\overline{k}} \cap \overline{V}$ sont quasi-compacts et, d'après le
lemme \ref{lemma-Galois-action-quasi-compact-open},
$W_{\overline{k}} \cap \overline{U}$ et
$W_{\overline{k}} \cap \overline{V}$
sont tous deux définis sur une sous-extension finie et stables par un
sous-groupe ouvert de $\text{Gal}(\overline{k}/k)$.
Nous l'utiliserons désormais sans autre mention.

\medskip\noindent
Choisissons un sous-espace ouvert quasi-compact $W_0 \subset X$ tel que
$W_{0, \overline{k}} \cap \overline{U}$ et
$W_{0, \overline{k}} \cap \overline{V}$ soient tous deux non vides.
Choisissons une sous-extension finie $\overline{k}/k'/k$
et une décomposition $W_{0, k'} = U_0' \amalg V_0'$ en parties ouvertes et fermées
telle que
$W_{0, \overline{k}} \cap \overline{U} = (U'_0)_{\overline{k}}$ et
$W_{0, \overline{k}} \cap \overline{V} = (V'_0)_{\overline{k}}$.
Posons $H = \text{Gal}(\overline{k}/k') \subset \text{Gal}(\overline{k}/k)$.
En particulier,
$\sigma(W_{0, \overline{k}} \cap \overline{U}) =
W_{0, \overline{k}} \cap \overline{U}$, et de même pour
$\overline{V}$.

\medskip\noindent
$W_0$, $k'$ étant choisis comme ci-dessus, pour tout sous-espace ouvert quasi-compact
$W \subset X$, posons
$$
U_W =
\bigcap\nolimits_{\sigma \in H} \sigma(W_{\overline{k}} \cap \overline{U}),
\quad
V_W =
\bigcup\nolimits_{\sigma \in H} \sigma(W_{\overline{k}} \cap \overline{V}).
$$
Comme $W_{\overline{k}} \cap \overline{U}$ et
$W_{\overline{k}} \cap \overline{V}$ sont fixés par un sous-groupe ouvert de
$\text{Gal}(\overline{k}/k)$, la réunion et l'intersection
ci-dessus sont finies. Ainsi, $U_W$ et $V_W$ sont des sous-espaces ouverts et fermés.
De plus, par construction, $W_{\bar k} = U_W \amalg V_W$.

\medskip\noindent
Affirmons que, si $W \subset W' \subset X$ sont des sous-espaces ouverts
quasi-compacts, alors $W_{\overline{k}} \cap U_{W'} = U_W$ et
$W_{\overline{k}} \cap V_{W'} = V_W$. Vérification omise.
Il s'ensuit qu'en posant $U = \bigcup_{W \subset X} U_W$
et $V = \bigcup_{W \subset X} V_W$, nous obtenons
$X_{\overline{k}} = U \amalg V$ comme réunion disjointe de parties ouvertes
et fermées.
Il est clair que $V$ est non vide, puisqu'il est construit en prenant des
réunions (localement). D'autre part, $U$ est non vide puisqu'il contient
$W_0 \cap \overline{U}$ par construction. Enfin, $U, V \subset X_{\bar k}$
sont fermés et stables par $H$ par construction. Donc, d'après le
lemme \ref{lemma-closed-fixed-by-Galois},
nous avons $U = (U')_{\bar k}$ et $V = (V')_{\bar k}$ pour certaines
parties fermées $U', V' \subset X_{k'}$. Manifestement, $X_{k'} = U' \amalg V'$,
et $X_{k'}$ est non connexe, comme voulu.
\end{proof}






\section{Espaces algébriques géométriquement irréductibles}
\label{section-geometrically-irreducible}

\noindent
L'exemple d'Espaces \ref{spaces-example-infinite-product} montre qu'il
vaut mieux ne pas considérer les espaces algébriques irréductibles en toute
généralité\footnote{Bien entendu, si nous disons « l'espace algébrique $X$
est irréductible », nous voulons probablement dire « l'espace topologique $|X|$
est irréductible ».}. Pour les espaces algébriques décents (par exemple quasi-séparés),
ce type de désastre ne se produit pas. Nous posons donc la
définition suivante seulement sous l'hypothèse que notre espace algébrique est décent.

\begin{definition}
\label{definition-geometrically-irreducible}
Soit $k$ un corps.
Soit $X$ un espace algébrique décent sur $k$.
Nous disons que $X$ est {\it géométriquement irréductible} si l'espace
topologique $|X_{k'}|$ est
irréductible\footnote{Un espace irréductible est non vide.}
pour toute extension de corps $k'$ de $k$.
\end{definition}

\noindent
Observons que $X_{k'}$ est un espace algébrique décent
(Espaces décents, lemme
\ref{decent-spaces-lemma-representable-named-properties}).
Par conséquent, l'espace topologique $|X_{k'}|$ est sobre,
d'après la proposition d'Espaces décents \ref{decent-spaces-proposition-reasonable-sober}.





\section{Espaces algébriques géométriquement intègres}
\label{section-geometrically-integral}

\noindent
Rappelons que les espaces algébriques intègres sont décents par définition ; voir
la section \ref{section-integral-spaces}.

\begin{definition}
\label{definition-geometrically-integral}
Soit $X$ un espace algébrique sur le corps $k$. Nous disons que $X$ est
{\it géométriquement intègre} sur $k$ si l'espace algébrique
$X_{k'}$ est intègre (définition \ref{definition-integral-algebraic-space})
pour toute extension de corps $k'$ de $k$.
\end{definition}

\noindent
En particulier, $X$ est un espace algébrique décent.
On peut relier cette propriété au fait d'être géométriquement réduit et
géométriquement irréductible comme suit.

\begin{lemma}
\label{lemma-geometrically-integral}
Soit $k$ un corps. Soit $X$ un espace algébrique décent sur $k$.
Alors $X$ est géométriquement intègre sur $k$ si et seulement si
$X$ est à la fois géométriquement réduit et géométriquement irréductible
sur $k$.
\end{lemma}

\begin{proof}
C'est une conséquence immédiate des définitions, car notre
notion d'intégrité (sous l'hypothèse de décence) équivaut à
être réduit et irréductible.
\end{proof}

\begin{lemma}
\label{lemma-proper-geometrically-reduced-global-sections}
Soit $k$ un corps. Soit $X$ un espace algébrique propre sur $k$.
\begin{enumerate}
\item $A = H^0(X, \mathcal{O}_X)$ est une $k$-algèbre de dimension finie ;
\item $A = \prod_{i = 1, \ldots, n} A_i$ est un produit de
$k$-algèbres locales artiniennes, avec un facteur par composante connexe de $|X|$ ;
\item si $X$ est réduit, alors $A = \prod_{i = 1, \ldots, n} k_i$
est un produit de corps, chacun étant une extension finie de $k$ ;
\item si $X$ est géométriquement réduit, alors $k_i$ est fini séparable
sur $k$ ;
\item si $X$ est géométriquement connexe, alors $A$ est géométriquement
irréductible sur $k$ ;
\item si $X$ est géométriquement irréductible, alors $A$ est géométriquement
irréductible sur $k$ ;
\item si $X$ est géométriquement réduit et connexe, alors $A = k$ ;
\item si $X$ est géométriquement intègre, alors $A = k$.
\end{enumerate}
\end{lemma}

\begin{proof}
D'après le lemme de Cohomologie des espaces
\ref{spaces-cohomology-lemma-proper-over-affine-cohomology-finite}
$A = H^0(X, \mathcal{O}_X)$ est une
$k$-algèbre de dimension finie. Cela démontre (1).

\medskip\noindent
Alors $A$ est un produit d'anneaux locaux d'après le
lemme d'Algèbre \ref{algebra-lemma-finite-dimensional-algebra} et la
proposition d'Algèbre \ref{algebra-proposition-dimension-zero-ring}.
Si $X = Y \amalg Z$, avec $Y$ et $Z$ des sous-espaces ouverts de $X$, nous obtenons
un idempotent $e \in A$ en prenant la section de $\mathcal{O}_X$
qui vaut $1$ sur $Y$ et $0$ sur $Z$. Réciproquement, si $e \in A$
est un idempotent, nous obtenons une décomposition correspondante de $|X|$.
Enfin, comme $|X|$ est un espace topologique noethérien
(d'après le lemme de Morphismes d'espaces
\ref{spaces-morphisms-lemma-finite-presentation-noetherian} et
le lemme de Propriétés des espaces
\ref{spaces-properties-lemma-Noetherian-topology})
ses composantes connexes sont ouvertes. Les composantes connexes
de $|X|$ correspondent donc $1$ pour $1$ aux idempotents primitifs de $A$.
Cela démontre (2).

\medskip\noindent
Si $X$ est réduit, alors $A$ est réduit
(Propriétés des espaces, lemme \ref{spaces-properties-lemma-reduced-space}).
Les anneaux locaux $A_i = k_i$ sont donc réduits et sont par conséquent des corps
(par exemple d'après le lemme d'Algèbre \ref{algebra-lemma-minimal-prime-reduced-ring}).
Cela démontre (3).

\medskip\noindent
Si $X$ est géométriquement réduit, la même propriété vaut pour
$A \otimes_k \overline{k} =
H^0(X_{\overline{k}}, \mathcal{O}_{X_{\overline{k}}})$
(voir le lemme de Cohomologie des espaces
\ref{spaces-cohomology-lemma-flat-base-change-cohomology} pour l'égalité).
Il s'ensuit que $k_i \otimes_k \overline{k}$ est un produit
de corps et que $k_i/k$ est donc séparable, par exemple d'après les
lemmes d'Algèbre
\ref{algebra-lemma-characterize-separable-field-extensions} et
\ref{algebra-lemma-geometrically-reduced-finite-purely-inseparable-extension}.
Cela démontre (4).

\medskip\noindent
Si $X$ est géométriquement connexe, alors $A \otimes_k \overline{k} =
H^0(X_{\overline{k}}, \mathcal{O}_{X_{\overline{k}}})$
est un anneau local de dimension zéro d'après (2) ; son
spectre n'a donc qu'un point et il est en particulier irréductible.
Ainsi, $A$ est géométriquement irréductible. Cela démontre (5).
Bien entendu, (5) entraîne (6).

\medskip\noindent
Si $X$ est géométriquement réduit et connexe, alors
$A = k_1$ est un corps et l'extension $k_1/k$ est finie séparable et
géométriquement irréductible. Mais $k_1 \otimes_k \overline{k}$
est alors un produit de $[k_1 : k]$ copies de $\overline{k}$, d'où
$k_1 = k$. Cela démontre (7). Bien entendu, (7) entraîne (8).
\end{proof}

\begin{lemma}
\label{lemma-characterize-trivial-pic-integral}
Soit $k$ un corps. Soit $X$ un espace algébrique propre intègre sur $k$.
Soit $\mathcal{L}$ un $\mathcal{O}_X$-module inversible.
Si $H^0(X, \mathcal{L})$ et $H^0(X, \mathcal{L}^{\otimes - 1})$
sont tous deux non nuls, alors $\mathcal{L} \cong \mathcal{O}_X$.
\end{lemma}

\begin{proof}
Soient $s \in H^0(X, \mathcal{L})$ et $t \in H^0(X, \mathcal{L}^{\otimes - 1})$
des sections non nulles. Soit $x \in |X|$ un point appartenant au support de $s$.
Choisissons un voisinage affine \'etale $(U, u) \to (X, x)$ tel que
$\mathcal{L}|_U \cong \mathcal{O}_U$. Alors $s|_U$ correspond à une fonction régulière non nulle
sur un schéma réduit (car $X$ est réduit), à savoir $U$, et elle ne
s'annule donc pas en un point générique d'une composante irréductible de $U$. D'après le
lemme d'Espaces décents \ref{decent-spaces-lemma-decent-generic-points},
le point générique $\eta$ de $|X|$ appartient au support de $s$.
Il en va de même pour $t$. Alors, bien sûr, $st$ est non nul, car
l'anneau local de $X$ en $\eta$ est un corps (d'après le lemme précité,
l'anneau local est de dimension zéro ; comme $X$ est réduit, l'anneau local est
réduit ; appliquer le lemme d'Algèbre \ref{algebra-lemma-minimal-prime-reduced-ring}).
Or, nous avons vu que $K = H^0(X, \mathcal{O}_X)$
est un corps au lemme \ref{lemma-proper-geometrically-reduced-global-sections}.
Ainsi, $st$ est partout non nul, et
$s : \mathcal{O}_X \to \mathcal{L}$ est un isomorphisme.
\end{proof}






\section{Dimension}
\label{section-dimension}

\noindent
Dans cette section, nous poursuivons l'étude de la dimension.
Voici une liste des résultats antérieurs :
\begin{enumerate}
\item la dimension est définie dans la
section de Propriétés des espaces
\ref{spaces-properties-section-dimension},
\item la dimension de l'anneau local est définie dans la
section de Propriétés des espaces
\ref{spaces-properties-section-dimension-local-ring},
\item quelques résultats figurent dans les lemmes de Propriétés des espaces
\ref{spaces-properties-lemma-dimension-local-ring} et
\ref{spaces-properties-lemma-dimension-decent-invariant-under-etale},
\item la dimension relative est définie dans la
section de Morphismes d'espaces \ref{spaces-morphisms-section-relative-dimension},
\item des résultats sur la dimension des fibres figurent dans la
section de Morphismes d'espaces \ref{spaces-morphisms-section-dimension-fibres},
\item une forme faible de la formule de dimension figure dans la
section de Morphismes d'espaces \ref{spaces-morphisms-section-dimension-formula},
\item un résultat sur la lissité et la dimension figure dans le lemme de Morphismes d'espaces
\ref{spaces-morphisms-lemma-smoothness-dimension-spaces},
\item pour les espaces décents, la dimension est $\dim(|X|)$, d'après le
lemme d'Espaces décents \ref{decent-spaces-lemma-dimension-decent-space},
\item les morphismes quasi-finis et la dimension sont étudiés dans les
lemmes d'Espaces décents
\ref{decent-spaces-lemma-dimension-local-ring-quasi-finite} et
\ref{decent-spaces-lemma-dimension-quasi-finite}.
\end{enumerate}
Dans la section de Compléments sur les morphismes d'espaces
\ref{spaces-more-morphisms-section-dimension}
nous étudierons les sauts de
dimension dans les fibres d'un morphisme de type fini.

\begin{lemma}
\label{lemma-integral-dimension}
Soit $S$ un schéma. Soit $f : X \to Y$ un morphisme entier
d'espaces algébriques. Alors $\dim(X) \leq \dim(Y)$.
Si $f$ est surjectif, alors $\dim(X) = \dim(Y)$.
\end{lemma}

\begin{proof}
Choisissons $V \to Y$ \'etale surjectif, avec $V$ un schéma.
Alors $U = X \times_Y V$ est un schéma et $U \to V$ est entier
(et surjectif si $f$ est surjectif).
D'après le lemme de Propriétés des espaces
\ref{spaces-properties-lemma-dimension-decent-invariant-under-etale}
nous avons $\dim(X) = \dim(U)$ et $\dim(Y) = \dim(V)$.
Le résultat découle donc du cas des schémas,
qui est le lemme de Morphismes \ref{morphisms-lemma-integral-dimension}.
\end{proof}

\begin{lemma}
\label{lemma-alteration-dimension}
Soit $S$ un schéma. Soit $f : X \to Y$ un morphisme d'espaces algébriques
sur $S$. Supposons que :
\begin{enumerate}
\item $Y$ est localement noethérien ;
\item $X$ et $Y$ sont des espaces algébriques intègres ;
\item $f$ est dominant ;
\item $f$ est localement de type fini.
\end{enumerate}
Si $x \in |X|$ et $y \in |Y|$ sont les points génériques, alors
$$
\dim(X) \leq \dim(Y) + \text{degré de transcendance de }x/y.
$$
Si $f$ est propre, alors il y a égalité.
\end{lemma}

\begin{proof}
Rappelons que $|X|$ et $|Y|$ sont des espaces topologiques irréductibles et sobres ; voir la
discussion qui suit la définition \ref{definition-integral-algebraic-space}.
Ainsi, la dominance de $f$ signifie que $|f|$ envoie $x$ sur $y$.
De plus, $x \in |X|$ est l'unique point auquel
l'anneau local de $X$ est de dimension $0$ ; voir le
lemme d'Espaces décents \ref{decent-spaces-lemma-decent-generic-points}.
D'après le lemme de Morphismes d'espaces
\ref{spaces-morphisms-lemma-dimension-formula-general}
la dimension de l'anneau local de $X$ en
tout point $x' \in |X|$ est au plus égale à celle de l'anneau local
de $Y$ en $y' = f(x')$, augmentée du degré de transcendance de $x/y$.
Comme la dimension de $X$, resp.\ celle de $Y$, est le
supremum des dimensions des anneaux locaux en $x'$, resp.\ en $y'$
(Propriétés des espaces, lemme \ref{spaces-properties-lemma-dimension}),
nous obtenons l'inégalité.

\medskip\noindent
Supposons $f$ propre.
Soit $V \subset Y$ un sous-espace ouvert quasi-compact non vide.
Si nous pouvons démontrer l'égalité pour le morphisme $f^{-1}(V) \to V$,
nous l'obtenons pour $X \to Y$. Nous pouvons donc supposer que
$X$ et $Y$ sont quasi-compacts.
Observons que $X$ est quasi-séparé, car c'est
un espace algébrique décent localement noethérien ; voir le
lemme d'Espaces décents
\ref{decent-spaces-lemma-locally-Noetherian-decent-quasi-separated}.
Nous pouvons donc choisir $Y' \to Y$ fini surjectif, où $Y'$
est un schéma ; voir la proposition de Limites d'espaces
\ref{spaces-limits-proposition-there-is-a-scheme-finite-over}.
Après avoir remplacé $Y'$ par un sous-schéma fermé convenable, nous
pouvons supposer $Y'$ intègre ; voir, par exemple, le lemme plus général
\ref{lemma-alteration-contained-in}.
D'après le même lemme, nous pouvons choisir un sous-espace fermé
$X' \subset X \times_Y Y'$ tel que $X'$ soit intègre
et que $X' \to X$ soit fini surjectif.
Alors $X'$ est aussi localement noethérien
(Morphismes d'espaces, lemme
\ref{spaces-morphisms-lemma-locally-finite-type-locally-noetherian})
et nous pouvons utiliser la proposition de Limites d'espaces
\ref{spaces-limits-proposition-there-is-a-scheme-finite-over}
une nouvelle fois pour choisir un morphisme fini surjectif $X'' \to X'$
avec $X''$ un schéma. Comme précédemment, nous pouvons supposer $X''$
intègre. On obtient le diagramme
$$
\xymatrix{
X'' \ar[d] \ar[r] & X \ar[d]^f \\
Y' \ar[r] & Y
}
$$
D'après le lemme \ref{lemma-integral-dimension}, nous avons $\dim(X'') = \dim(X)$
et $\dim(Y') = \dim(Y)$. Comme $X$ et $Y$ possèdent des voisinages ouverts
de $x$, resp.\ de $y$, qui sont des schémas, nous voyons aussitôt que les points génériques
$x'' \in X''$, resp.\ $y' \in Y'$, sont les uniques points envoyés
sur $x$, resp.\ sur $y$, et que les extensions de corps résiduels
$\kappa(x'')/\kappa(x)$ et $\kappa(y')/\kappa(y)$ sont finies.
Il s'ensuit que le degré de transcendance de $x''/y'$ est le
même que celui de $x/y$. L'égalité découle donc
du cas des schémas, qui est le
lemme de Morphismes \ref{morphisms-lemma-alteration-dimension}.
\end{proof}





\section{Espaces lisses sur les corps}
\label{section-smooth}

\noindent
Cette section est l'analogue de la
section de Variétés \ref{varieties-section-smooth}.

\begin{lemma}
\label{lemma-smooth-regular}
Soit $k$ un corps.
Soit $X$ un espace algébrique lisse sur $k$.
Alors $X$ est un espace algébrique régulier.
\end{lemma}

\begin{proof}
Choisissons un schéma $U$ et un morphisme \'etale surjectif $U \to X$.
Le morphisme $U \to \Spec(k)$ est lisse, car il est composé
d'un morphisme \'etale (donc lisse) et d'un morphisme lisse (voir les
lemmes de Morphismes d'espaces \ref{spaces-morphisms-lemma-etale-smooth}
et \ref{spaces-morphisms-lemma-composition-smooth}).
Par conséquent, $U$ est régulier d'après le
lemme de Variétés \ref{varieties-lemma-smooth-regular}.
D'après la
définition de Propriétés des espaces
\ref{spaces-properties-definition-type-property}
cela signifie que $X$ est régulier.
\end{proof}

\begin{lemma}
\label{lemma-smooth-separable-closed-points-dense}
Soit $k$ un corps. Soit $X$ un espace algébrique lisse sur $\Spec(k)$.
L'ensemble des $x \in |X|$ qui sont images de morphismes $\Spec(k') \to X$
où $k' \supset k$ est fini séparable est dense dans $|X|$.
\end{lemma}

\begin{proof}
Choisissons un schéma $U$ et un morphisme \'etale surjectif $U \to X$.
Le morphisme $U \to \Spec(k)$ est lisse, car il est composé
d'un morphisme \'etale (donc lisse) et d'un morphisme lisse (voir les
lemmes de Morphismes d'espaces \ref{spaces-morphisms-lemma-etale-smooth}
et \ref{spaces-morphisms-lemma-composition-smooth}).
Nous pouvons donc appliquer le lemme de Variétés
\ref{varieties-lemma-smooth-separable-closed-points-dense}, selon lequel
les points fermés de $U$ dont les corps résiduels sont finis séparables sur
$k$ sont denses. Le lemme en résulte par notre définition de la
topologie sur $|X|$.
\end{proof}







\section{Caractéristiques d'Euler-Poincaré}
\label{section-euler}

\noindent
Dans cette section, nous démontrons quelques propriétés élémentaires des caractéristiques
d'Euler-Poincaré des faisceaux cohérents sur les espaces algébriques propres sur des corps.

\begin{definition}
\label{definition-euler-characteristic}
Soit $k$ un corps. Soit $X$ un espace algébrique propre sur $k$. Soit $\mathcal{F}$
un $\mathcal{O}_X$-module cohérent. Dans cette situation, la
{\it caractéristique d'Euler-Poincaré de $\mathcal{F}$} est l'entier
$$
\chi(X, \mathcal{F}) = \sum\nolimits_i (-1)^i \dim_k H^i(X, \mathcal{F}).
$$
La justification de cette formule est donnée ci-dessous.
\end{definition}

\noindent
Dans la situation de la définition, les espaces vectoriels
$H^i(X, \mathcal{F})$ non nuls sont en nombre fini (lemme de Cohomologie des espaces
\ref{spaces-cohomology-lemma-vanishing-quasi-separated})
et ils sont tous de dimension finie
(lemme de Cohomologie des espaces
\ref{spaces-cohomology-lemma-proper-over-affine-cohomology-finite}). Ainsi,
$\chi(X, \mathcal{F}) \in \mathbf{Z}$ est bien défini. Remarquons que
cette définition dépend du corps $k$ et non pas seulement de la paire
$(X, \mathcal{F})$.

\begin{lemma}
\label{lemma-euler-characteristic-additive}
Soit $k$ un corps. Soit $X$ un espace algébrique propre sur $k$.
Soit $0 \to \mathcal{F}_1 \to \mathcal{F}_2 \to \mathcal{F}_3 \to 0$
une suite exacte courte de modules cohérents sur $X$. Alors
$$
\chi(X, \mathcal{F}_2) = \chi(X, \mathcal{F}_1) + \chi(X, \mathcal{F}_3)
$$
\end{lemma}

\begin{proof}
Considérons la suite exacte longue de cohomologie
$$
0 \to H^0(X, \mathcal{F}_1) \to H^0(X, \mathcal{F}_2) \to
H^0(X, \mathcal{F}_3) \to H^1(X, \mathcal{F}_1) \to \ldots
$$
associée à la suite exacte courte du lemme. Le théorème du rang
en algèbre linéaire montre que
$$
0 = \dim H^0(X, \mathcal{F}_1) - \dim H^0(X, \mathcal{F}_2)
+ \dim H^0(X, \mathcal{F}_3) - \dim H^1(X, \mathcal{F}_1) + \ldots
$$
Cela implique immédiatement le lemme.
\end{proof}

\begin{lemma}
\label{lemma-euler-characteristic-morphism}
Soit $k$ un corps. Soit $f : Y \to X$ un morphisme d'
espaces algébriques propres sur $k$. Soit $\mathcal{G}$ un
$\mathcal{O}_Y$-module cohérent. Alors
$$
\chi(Y, \mathcal{G}) = \sum (-1)^i \chi(X, R^if_*\mathcal{G})
$$
\end{lemma}

\begin{proof}
La formule a un sens : les faisceaux $R^if_*\mathcal{G}$ sont cohérents
et ceux qui sont non nuls sont en nombre fini ; voir les
lemmes de Cohomologie des espaces
\ref{spaces-cohomology-lemma-proper-pushforward-coherent} et
\ref{spaces-cohomology-lemma-vanishing-higher-direct-images}.
D'après le lemme de Cohomologie sur les sites \ref{sites-cohomology-lemma-Leray},
il existe une suite spectrale telle que
$$
E_2^{p, q} = H^p(X, R^qf_*\mathcal{G})
$$
qui aboutit à $H^{p + q}(Y, \mathcal{G})$. La finitude de la cohomologie
sur $X$ montre que les $E_2^{p, q}$ non nuls sont en nombre fini
et que chaque $E_2^{p, q}$ est un espace vectoriel de dimension finie. Il s'ensuit
qu'il en va de même pour $E_r^{p, q}$ lorsque $r \geq 2$ et que
$$
\sum (-1)^{p + q} \dim_k E_r^{p, q}
$$
est indépendante de $r$. Comme, pour $r$ assez grand, nous avons
$E_r^{p, q} = E_\infty^{p, q}$ et que la convergence signifie qu'il
existe sur $H^n(Y, \mathcal{G})$ une filtration dont les quotients gradués sont
les $E_\infty^{p, q}$ avec $p + 1 = n$ (c'est le sens de la convergence
de la suite spectrale), nous concluons.
\end{proof}








\section{Nombres d'intersection}
\label{section-num}

\noindent
Dans cette section, nous étudions la caractéristique d'Euler-Poincaré des
faisceaux cohérents sur les espaces algébriques propres afin d'obtenir des
nombres d'intersection pour les modules inversibles. Notre outil principal sera le
lemme suivant.

\begin{lemma}
\label{lemma-numerical-polynomial-from-euler}
Soit $k$ un corps. Soit $X$ un espace algébrique propre sur $k$.
Soit $\mathcal{F}$ un $\mathcal{O}_X$-module cohérent. Soient
$\mathcal{L}_1, \ldots, \mathcal{L}_r$ des $\mathcal{O}_X$-modules inversibles.
L'application
$$
(n_1, \ldots, n_r) \longmapsto
\chi(X, \mathcal{F} \otimes
\mathcal{L}_1^{\otimes n_1} \otimes \ldots \otimes
\mathcal{L}_r^{\otimes n_r})
$$
est un polynôme numérique en $n_1, \ldots, n_r$ de degré total au
plus égal à la dimension du support schématique de $\mathcal{F}$.
\end{lemma}

\begin{proof}
Soit $Z \subset X$ le support schématique de $\mathcal{F}$.
Alors $\mathcal{F} = i_*\mathcal{G}$ pour un
$\mathcal{O}_Z$-module cohérent $\mathcal{G}$
(lemme de Cohomologie des espaces
\ref{spaces-cohomology-lemma-coherent-support-closed})
et nous avons
$$
\chi(X, \mathcal{F} \otimes
\mathcal{L}_1^{\otimes n_1} \otimes \ldots \otimes
\mathcal{L}_r^{\otimes n_r}) =
\chi(Z, \mathcal{G} \otimes
i^*\mathcal{L}_1^{\otimes n_1} \otimes \ldots \otimes
i^*\mathcal{L}_r^{\otimes n_r})
$$
d'après la formule de projection
(lemme de Cohomologie sur les sites \ref{sites-cohomology-lemma-projection-formula})
et le lemme de Cohomologie des espaces
\ref{spaces-cohomology-lemma-relative-affine-cohomology}.
Comme $|Z| = \text{Supp}(\mathcal{F})$, il suffit
de montrer que
$$
P_\mathcal{F}(n_1, \ldots, n_r) :
(n_1, \ldots, n_r)
\longmapsto
\chi(X, \mathcal{F} \otimes
\mathcal{L}_1^{\otimes n_1} \otimes \ldots \otimes
\mathcal{L}_r^{\otimes n_r})
$$
est un polynôme numérique en $n_1, \ldots, n_r$ de degré total au
plus égal à $\dim(X)$. Disons que la propriété $\mathcal{P}$ est vérifiée par le
$\mathcal{O}_X$-module cohérent $\mathcal{F}$ si l'assertion ci-dessus est vraie.

\medskip\noindent
Nous démontrerons cet énoncé par dévissage ; plus précisément, nous
vérifierons que les conditions (1), (2) et (3) du
lemme de Cohomologie des espaces
\ref{spaces-cohomology-lemma-property-higher-rank-cohomological-variant}
sont satisfaites.

\medskip\noindent
Vérification de la condition (1). Soit
$$
0 \to \mathcal{F}_1 \to \mathcal{F}_2 \to \mathcal{F}_3 \to 0
$$
une suite exacte courte de faisceaux cohérents sur $X$.
D'après le lemme \ref{lemma-euler-characteristic-additive}, nous avons
$$
P_{\mathcal{F}_2}(n_1, \ldots, n_r) =
P_{\mathcal{F}_1}(n_1, \ldots, n_r) +
P_{\mathcal{F}_3}(n_1, \ldots, n_r)
$$
Il est alors clair que, si deux des trois faisceaux $\mathcal{F}_i$
vérifient la propriété $\mathcal{P}$, le troisième la vérifie aussi.

\medskip\noindent
La condition (2) résulte de l'égalité
$P_{\mathcal{F}^{\oplus m}}(n_1, \ldots, n_r) =
mP_\mathcal{F}(n_1, \ldots, n_r)$.

\medskip\noindent
Démonstration de (3). Soit $i : Z \to X$ un sous-espace fermé réduit tel que
$|Z|$ soit irréductible. Nous devons trouver un module cohérent $\mathcal{G}$
sur $X$, de support $Z$, tel que $\mathcal{P}$ soit vérifiée pour $\mathcal{G}$.
Nous donnerons deux constructions : l'une au moyen du lemme de Chow, l'autre
au moyen d'un revêtement fini par un schéma.

\medskip\noindent
Existence de $\mathcal{G}$ au moyen d'un revêtement fini par un schéma.
Choisissons $\pi : Z' \to Z$ fini surjectif, où $Z'$ est un schéma ; voir la
proposition de Limites d'espaces
\ref{spaces-limits-proposition-there-is-a-scheme-finite-over}.
Posons $\mathcal{G} = i_*\pi_*\mathcal{O}_{Z'} = (i \circ \pi)_*\mathcal{O}_{Z'}$.
Remarquons que $Z'$ est propre sur $k$ et que le support de $\mathcal{G}$ est $Y$
(détails omis). Nous avons
$$
R(\pi \circ i)_*(\mathcal{O}_{Z'}) = \mathcal{G}
\quad\text{et}\quad
R(\pi \circ i)_*(\pi^*i^*(\mathcal{L}_1^{\otimes n_1} \otimes \ldots \otimes
\mathcal{L}_r^{\otimes n_r})
) = \mathcal{G} \otimes \mathcal{L}_1^{\otimes n_1} \otimes \ldots \otimes
\mathcal{L}_r^{\otimes n_r}
$$
La première égalité vaut parce que $i \circ \pi$ est affine
(lemme de Cohomologie des espaces
\ref{spaces-cohomology-lemma-affine-vanishing-higher-direct-images})
et la seconde résulte de la première et de la formule de projection
(lemme de Cohomologie sur les sites \ref{sites-cohomology-lemma-projection-formula}).
En utilisant Leray
(lemme de Cohomologie sur les sites \ref{sites-cohomology-lemma-apply-Leray}),
nous obtenons
$$
P_\mathcal{G}(n_1, \ldots, n_r) =
\chi(Z', \pi^*i^*(\mathcal{L}_1^{\otimes n_1} \otimes \ldots \otimes
\mathcal{L}_r^{\otimes n_r}))
$$
D'après le cas des schémas
(lemme de Variétés \ref{varieties-lemma-numerical-polynomial-from-euler}),
il s'agit d'un polynôme numérique en
$n_1, \ldots, n_r$ de degré au plus égal à $\dim(Z')$.
Nous concluons puisque $\dim(Z') \leq \dim(Z) \leq \dim(X)$.
La première inégalité résulte du
lemme d'Espaces décents \ref{decent-spaces-lemma-dimension-quasi-finite}.

\medskip\noindent
Existence de $\mathcal{G}$ au moyen du lemme de Chow. Nous appliquons le
lemme de Cohomologie des espaces \ref{spaces-cohomology-lemma-weak-chow}
au morphisme $Z \to \Spec(k)$. Nous obtenons ainsi un
morphisme propre surjectif $f : Y \to Z$ sur $\Spec(k)$,
où $Y$ est un sous-schéma fermé de $\mathbf{P}^m_k$ pour un certain $m$.
Après remplacement de $Y$ par un sous-schéma fermé, nous pouvons supposer que $Y$
est intègre et que $f : Y \to Z$ est une altération ; voir le
lemme \ref{lemma-alteration-contained-in}.
Notons $\mathcal{O}_Y(n)$ le tiré en arrière de $\mathcal{O}_{\mathbf{P}^m_k}(n)$.
Choisissons $n > 0$ tel que $R^pf_*\mathcal{O}_Y(n) = 0$
pour $p > 0$ ; voir le lemme de Cohomologie des espaces
\ref{spaces-cohomology-lemma-kill-by-twisting}.
Nous affirmons que $\mathcal{G} = i_*f_*\mathcal{O}_Y(n)$ vérifie $\mathcal{P}$.
En effet, le cas des schémas
(lemme de Variétés \ref{varieties-lemma-numerical-polynomial-from-euler})
montre que
$$
(n_1, \ldots, n_r)
\longmapsto
\chi(Y, \mathcal{O}_Y(n) \otimes
f^*i^*(\mathcal{L}_1^{\otimes n_1} \otimes \ldots \otimes
\mathcal{L}_r^{\otimes n_r}))
$$
est un polynôme numérique en $n_1, \ldots, n_r$ de degré total au
plus égal à $\dim(Y)$. D'autre part, d'après la formule de projection
(lemme de Cohomologie sur les sites \ref{sites-cohomology-lemma-projection-formula}),
\begin{align*}
i_*Rf_*\left(
\mathcal{O}_Y(n) \otimes
f^*i^*(\mathcal{L}_1^{\otimes n_1} \otimes \ldots \otimes
\mathcal{L}_r^{\otimes n_r})\right)
& =
i_*Rf_*\mathcal{O}_Y(n) \otimes
\mathcal{L}_1^{\otimes n_1} \otimes \ldots \otimes
\mathcal{L}_r^{\otimes n_r} \\
& =
\mathcal{G} \otimes \mathcal{L}_1^{\otimes n_1} \otimes \ldots \otimes
\mathcal{L}_r^{\otimes n_r}
\end{align*}
la dernière égalité résultant de notre choix de $n$. D'après
Leray (lemme de Cohomologie sur les sites \ref{sites-cohomology-lemma-apply-Leray}),
nous obtenons
$$
\chi(Y, \mathcal{O}_Y(n) \otimes
f^*i^*(\mathcal{L}_1^{\otimes n_1} \otimes \ldots \otimes
\mathcal{L}_r^{\otimes n_r})) =
P_\mathcal{G}(n_1, \ldots, n_r)
$$
et nous concluons puisque $\dim(Y) \leq \dim(Z) \leq \dim(X)$.
La première inégalité résulte du
lemme de Morphismes d'espaces
\ref{spaces-morphisms-lemma-alteration-dimension-general}
et du fait que $Y \to Z$ est une altération (de sorte que
l'extension induite des corps résiduels aux points génériques est finie).
\end{proof}

\noindent
Le lemme suivant montre, en gros, que le coefficient dominant ne dépend que
de la longueur du module cohérent aux points génériques de son
support.

\begin{lemma}
\label{lemma-numerical-polynomial-leading-term}
Soit $k$ un corps. Soit $X$ un espace algébrique propre sur $k$. Soit
$\mathcal{F}$ un $\mathcal{O}_X$-module cohérent. Soient
$\mathcal{L}_1, \ldots, \mathcal{L}_r$ des $\mathcal{O}_X$-modules inversibles.
Posons $d = \dim(\text{Supp}(\mathcal{F}))$.
Soient $Z_i \subset X$ les composantes irréductibles
de $\text{Supp}(\mathcal{F})$ de dimension $d$. Soit $\overline{x}_i$
un point générique géométrique de $Z_i$ et posons
$m_i = \text{longueur}_{\mathcal{O}_{X, \overline{x}_i}}
(\mathcal{F}_{\overline{x}_i})$.
Alors
$$
\chi(X, \mathcal{F} \otimes \mathcal{L}_1^{\otimes n_1} \otimes \ldots \otimes
\mathcal{L}_r^{\otimes n_r}) -
\sum\nolimits_i
m_i\ \chi(Z_i, \mathcal{L}_1^{\otimes n_1} \otimes \ldots \otimes
\mathcal{L}_r^{\otimes n_r}|_{Z_i})
$$
est un polynôme numérique en $n_1, \ldots, n_r$ de degré total $< d$.
\end{lemma}

\begin{proof}
Démontrons d'abord un énoncé légèrement plus faible. Posons
$\dim(X) = N$ et soient $X_i \subset X$ les composantes
irréductibles de dimension $N$. Soit $\overline{x}_i$ un point
générique géométrique de $X_i$. L'anneau local \'etale
$\mathcal{O}_{X, \overline{x}_i}$ est noethérien de
dimension $0$ ; ainsi, pour tout $\mathcal{O}_X$-module cohérent
$\mathcal{F}$, la longueur
$$
m_i(\mathcal{F}) = \text{longueur}_{\mathcal{O}_{X, \overline{x}_i}}
(\mathcal{F}_{\overline{x}_i})
$$
est un entier $\geq 0$. Nous affirmons que
$$
E(\mathcal{F}) =
\chi(X, \mathcal{F} \otimes \mathcal{L}_1^{\otimes n_1} \otimes \ldots \otimes
\mathcal{L}_r^{\otimes n_r}) -
\sum\nolimits_i
m_i(\mathcal{F})\ \chi(Z_i, \mathcal{L}_1^{\otimes n_1} \otimes \ldots \otimes
\mathcal{L}_r^{\otimes n_r}|_{Z_i})
$$
est un polynôme numérique en $n_1, \ldots, n_r$ de degré total $< N$.
Nous le démontrerons au moyen du lemme de Cohomologie des espaces
\ref{spaces-cohomology-lemma-property-higher-rank-cohomological-variant}.
Pour toute suite exacte courte $0 \to \mathcal{F}' \to \mathcal{F} \to
\mathcal{F}'' \to 0$, nous avons
$E(\mathcal{F}) = E(\mathcal{F}') + E(\mathcal{F}'')$.
Cela résulte de l'additivité des caractéristiques d'Euler-Poincaré
(lemme \ref{lemma-euler-characteristic-additive})
et de l'additivité des longueurs
(lemme d'Algèbre \ref{algebra-lemma-length-additive}).
Cela implique immédiatement les propriétés (1) et (2) du lemme de Cohomologie des espaces
\ref{spaces-cohomology-lemma-property-higher-rank-cohomological-variant}.
Enfin, la propriété (3) est vérifiée avec $\mathcal{G} = \mathcal{O}_Z$
pour tout sous-espace fermé réduit et irréductible $Z \subset X$.
En effet, si $Z = Z_{i_0}$ pour un certain $i_0$, alors
$m_i(\mathcal{G}) = \delta_{i_0i}$ et nous concluons que $E(\mathcal{G}) = 0$.
Si $Z \not = Z_i$ pour tout $i$, alors $m_i(\mathcal{G}) = 0$ pour tout $i$,
$\dim(Z) < N$ et le résultat découle du
lemme \ref{lemma-numerical-polynomial-from-euler}.

\medskip\noindent
Démonstration de l'énoncé du lemme.
Soit $Z \subset X$ le support schématique de $\mathcal{F}$.
Alors $\mathcal{F} = i_*\mathcal{G}$ pour un
$\mathcal{O}_Z$-module cohérent $\mathcal{G}$
(lemme de Cohomologie des espaces
\ref{spaces-cohomology-lemma-coherent-support-closed})
et nous avons
$$
\chi(X, \mathcal{F} \otimes
\mathcal{L}_1^{\otimes n_1} \otimes \ldots \otimes
\mathcal{L}_r^{\otimes n_r}) =
\chi(Z, \mathcal{G} \otimes
i^*\mathcal{L}_1^{\otimes n_1} \otimes \ldots \otimes
i^*\mathcal{L}_r^{\otimes n_r})
$$
d'après la formule de projection
(lemme de Cohomologie sur les sites \ref{sites-cohomology-lemma-projection-formula})
et le lemme de Cohomologie des espaces
\ref{spaces-cohomology-lemma-relative-affine-cohomology}.
Comme $|Z| = \text{Supp}(\mathcal{F})$, nous avons $Z_i \subset Z$
pour tout $i$, et ce sont les composantes irréductibles
de $Z$ de dimension $d$. Nous considérons désormais $\overline{x}_i$
comme un point géométrique de $Z$. L'application
$i^\sharp : \mathcal{O}_X \to i_*\mathcal{O}_Z$
détermine une surjection
$$
\mathcal{O}_{X, \overline{x}_i} \to \mathcal{O}_{Z, \overline{x}_i}
$$
Cette application induit un isomorphisme de modules
$\mathcal{G}_{\overline{x}_i} = \mathcal{F}_{\overline{x}_i}$
puisque $\mathcal{F} = i_*\mathcal{G}$. Il en résulte que
$$
m_i = \text{longueur}_{\mathcal{O}_{X, \overline{x}_i}}
(\mathcal{F}_{\overline{x}_i}) =
\text{longueur}_{\mathcal{O}_{Z, \overline{x}_i}}
(\mathcal{G}_{\overline{x}_i})
$$
Nous voyons donc que l'expression du lemme est égale à
$$
\chi(Z, \mathcal{G} \otimes \mathcal{L}_1^{\otimes n_1} \otimes \ldots \otimes
\mathcal{L}_r^{\otimes n_r}) -
\sum\nolimits_i
m_i\ \chi(Z_i, \mathcal{L}_1^{\otimes n_1} \otimes \ldots \otimes
\mathcal{L}_r^{\otimes n_r}|_{Z_i})
$$
et le résultat découle de la discussion du premier paragraphe
(appliquée à $Z$ plutôt qu'à $X$).
\end{proof}

\begin{definition}
\label{definition-intersection-number}
Soit $k$ un corps. Soit $X$ un espace algébrique propre sur $k$. Soit
$i : Z \to X$ un sous-espace fermé de dimension $d$. Soient
$\mathcal{L}_1, \ldots, \mathcal{L}_d$ des
$\mathcal{O}_X$-modules inversibles. Nous définissons le {\it nombre d'intersection}
$(\mathcal{L}_1 \cdots \mathcal{L}_d \cdot Z)$
comme le coefficient de $n_1 \ldots n_d$ dans le polynôme numérique
$$
\chi(X, i_*\mathcal{O}_Z \otimes \mathcal{L}_1^{\otimes n_1} \otimes
\ldots \otimes \mathcal{L}_d^{\otimes n_d}) =
\chi(Z, \mathcal{L}_1^{\otimes n_1} \otimes
\ldots \otimes \mathcal{L}_d^{\otimes n_d}|_Z)
$$
Dans le cas particulier
où $\mathcal{L}_1 = \ldots = \mathcal{L}_d = \mathcal{L}$,
nous écrivons $(\mathcal{L}^d \cdot Z)$.
\end{definition}

\noindent
L'égalité affichée dans la définition résulte de
la formule de projection
(Cohomologie, section \ref{cohomology-section-projection-formula}) et du
lemme de Cohomologie des schémas
\ref{coherent-lemma-relative-affine-cohomology}.
Démontrons quelques lemmes sur ces nombres d'intersection.

\begin{lemma}
\label{lemma-intersection-number-integer}
Dans la situation de la définition \ref{definition-intersection-number},
le nombre d'intersection
$(\mathcal{L}_1 \cdots \mathcal{L}_d \cdot Z)$
est un entier.
\end{lemma}

\begin{proof}
Tout polynôme numérique de degré $e$ en $n_1, \ldots, n_d$ s'écrit
de manière unique comme combinaison $\mathbf{Z}$-linéaire des fonctions
${n_1 \choose k_1}{n_2 \choose k_2} \ldots {n_d \choose k_d}$ avec
$k_1 + \ldots + k_d \leq e$. Appliquons ceci avec $e = d$.
La vérification est laissée en exercice.
\end{proof}

\begin{lemma}
\label{lemma-intersection-number-additive}
Dans la situation de la définition \ref{definition-intersection-number},
le nombre d'intersection
$(\mathcal{L}_1 \cdots \mathcal{L}_d \cdot Z)$
est additif : si $\mathcal{L}_i = \mathcal{L}_i' \otimes \mathcal{L}_i''$,
alors nous avons
$$
(\mathcal{L}_1 \cdots \mathcal{L}_i \cdots \mathcal{L}_d \cdot Z) =
(\mathcal{L}_1 \cdots \mathcal{L}_i' \cdots \mathcal{L}_d \cdot Z) +
(\mathcal{L}_1 \cdots \mathcal{L}_i'' \cdots \mathcal{L}_d \cdot Z)
$$
\end{lemma}

\begin{proof}
Cela résulte du lemme \ref{lemma-numerical-polynomial-from-euler}, car
la fonction
$$
(n_1, \ldots, n_{i - 1}, n_i', n_i'', n_{i + 1}, \ldots, n_d)
\mapsto
\chi(Z, \mathcal{L}_1^{\otimes n_1} \otimes
\ldots \otimes (\mathcal{L}_i')^{\otimes n_i'} \otimes
(\mathcal{L}_i'')^{\otimes n_i''} \otimes \ldots \otimes
\mathcal{L}_d^{\otimes n_d}|_Z)
$$
est un polynôme numérique de degré total au plus égal à $d$ en $d + 1$ variables.
\end{proof}

\begin{lemma}
\label{lemma-intersection-number-in-terms-of-components}
Dans la situation de la définition \ref{definition-intersection-number},
soient $Z_i \subset Z$ les composantes irréductibles de dimension $d$. Soit
$m_i = \text{longueur}_{\mathcal{O}_{X, \overline{x}_i}}
(\mathcal{O}_{Z, \overline{x}_i})$
où $\overline{x}_i$ est un point générique géométrique de $Z_i$. Alors
$$
(\mathcal{L}_1 \cdots \mathcal{L}_d \cdot Z) =
\sum m_i(\mathcal{L}_1 \cdots \mathcal{L}_d \cdot Z_i)
$$
\end{lemma}

\begin{proof}
Cela résulte immédiatement du lemme \ref{lemma-numerical-polynomial-leading-term}
et des définitions.
\end{proof}

\begin{lemma}
\label{lemma-intersection-number-and-pullback}
Soit $k$ un corps. Soit $f : Y \to X$ un morphisme d'
espaces algébriques propres sur $k$.
Soit $Z \subset Y$ un sous-espace fermé intègre de dimension $d$ et soient
$\mathcal{L}_1, \ldots, \mathcal{L}_d$ des $\mathcal{O}_X$-modules inversibles.
Alors
$$
(f^*\mathcal{L}_1 \cdots f^*\mathcal{L}_d \cdot Z) =
\deg(f|_Z : Z \to f(Z)) (\mathcal{L}_1 \cdots \mathcal{L}_d \cdot f(Z))
$$
où $\deg(Z \to f(Z))$ est défini comme dans la définition \ref{definition-degree},
ou vaut $0$ si $\dim(f(Z)) < d$.
\end{lemma}

\begin{proof}
Dans l'énoncé, $f(Z) \subset X$ est l'image schématique de $f$
et c'est aussi l'espace algébrique muni de la structure réduite induite sur la
partie fermée $f(|Z|) \subset X$ ; voir le lemme de Morphismes d'espaces
\ref{spaces-morphisms-lemma-scheme-theoretic-image-reduced}.
Alors $Z$ et $f(Z)$ sont des espaces algébriques réduits, propres (donc décents)
sur $k$, et par suite intègres
(définition \ref{definition-integral-algebraic-space}).
Le membre de gauche se calcule comme le coefficient de $n_1 \ldots n_d$
dans la fonction
$$
\chi(Y, \mathcal{O}_Z \otimes f^*\mathcal{L}_1^{\otimes n_1} \otimes
\ldots \otimes f^*\mathcal{L}_d^{\otimes n_d}) =
\sum (-1)^i
\chi(X, R^if_*\mathcal{O}_Z \otimes
\mathcal{L}_1^{\otimes n_1} \otimes \ldots \otimes
\mathcal{L}_d^{\otimes n_d})
$$
L'égalité résulte du lemme \ref{lemma-euler-characteristic-morphism}
et de la formule de projection
(lemme de Cohomologie \ref{cohomology-lemma-projection-formula}).
Si $f(Z)$ est de dimension $< d$, alors le membre de droite
est un polynôme de degré total $< d$ d'après le
lemme \ref{lemma-numerical-polynomial-from-euler},
et le résultat est établi. Supposons $\dim(f(Z)) = d$. Alors
la théorie de la dimension (lemme \ref{lemma-alteration-dimension}) montre
que les conditions équivalentes (1) -- (5) du
lemme \ref{lemma-finite-degree} sont satisfaites. Ainsi,
$\deg(Z \to f(Z))$ est bien défini.
D'après le lemme \ref{lemma-finite-degree}, déjà utilisé,
$f : Z \to f(Z)$ est fini au-dessus d'un ouvert non vide
$V$ de $f(Z)$ ; après avoir éventuellement rétréci $V$, nous pouvons supposer que
$V$ est un schéma. Soit $\xi \in V$ le point générique.
Ainsi, $\deg(f : Z \to f(Z))$ est la longueur du germe de
$f_*\mathcal{O}_Z$ en $\xi$ sur $\mathcal{O}_{X, \xi}$,
et le germe de $R^if_*\mathcal{O}_X$ en $\xi$ est nul pour $i > 0$
(par exemple d'après le lemme de Cohomologie des espaces
\ref{spaces-cohomology-lemma-finite-higher-direct-image-zero}).
Ainsi, les termes $\chi(X, R^if_*\mathcal{O}_Z \otimes
\mathcal{L}_1^{\otimes n_1} \otimes \ldots \otimes
\mathcal{L}_d^{\otimes n_d})$ pour $i > 0$
ont un degré total $< d$ et
$$
\chi(X, f_*\mathcal{O}_Z \otimes
\mathcal{L}_1^{\otimes n_1} \otimes \ldots \otimes
\mathcal{L}_d^{\otimes n_d})
=
\deg(f : Z \to f(Z)) \chi(f(Z),
\mathcal{L}_1^{\otimes n_1} \otimes \ldots \otimes
\mathcal{L}_d^{\otimes n_d}|_{f(Z)})
$$
à un polynôme de degré total $< d$ près, d'après le
lemme \ref{lemma-numerical-polynomial-leading-term}.
Le résultat voulu en découle.
\end{proof}

\begin{lemma}
\label{lemma-numerical-intersection-effective-Cartier-divisor}
Soit $k$ un corps. Soit $X$ un espace algébrique propre sur $k$.
Soit $Z \subset X$ un sous-espace fermé de dimension $d$.
Soient $\mathcal{L}_1, \ldots, \mathcal{L}_d$
des $\mathcal{O}_X$-modules inversibles. Supposons qu'il existe un
diviseur de Cartier effectif $D \subset Z$ tel que
$\mathcal{L}_1|_Z \cong \mathcal{O}_Z(D)$. Alors
$$
(\mathcal{L}_1 \cdots \mathcal{L}_d \cdot Z) =
(\mathcal{L}_2 \cdots \mathcal{L}_d \cdot D)
$$
\end{lemma}

\begin{proof}
Nous pouvons remplacer $X$ par $Z$ et $\mathcal{L}_i$ par $\mathcal{L}_i|_Z$.
Ainsi, nous pouvons supposer $X = Z$ et $\mathcal{L}_1 = \mathcal{O}_X(D)$.
Alors $\mathcal{L}_1^{-1}$ est le faisceau d'idéaux de $D$ et nous pouvons
considérer la suite exacte courte
$$
0 \to \mathcal{L}_1^{\otimes -1} \to \mathcal{O}_X \to \mathcal{O}_D \to 0
$$
Posons
$P(n_1, \ldots, n_d) =
\chi(X, \mathcal{L}_1^{\otimes n_1} \otimes \ldots \otimes
\mathcal{L}_d^{\otimes n_d})$
et
$Q(n_1, \ldots, n_d) =
\chi(D, \mathcal{L}_1^{\otimes n_1} \otimes \ldots \otimes
\mathcal{L}_d^{\otimes n_d}|_D)$.
L'additivité
(lemme \ref{lemma-euler-characteristic-additive})
donne
$$
P(n_1, \ldots, n_d) - P(n_1 - 1, n_2, \ldots, n_d) =
Q(n_1, \ldots, n_d)
$$
Comme le degré total de $P$ est au plus égal à $d$, nous voyons que
le coefficient de $n_1 \ldots n_d$ dans $P$ est égal au coefficient
de $n_2 \ldots n_d$ dans $Q$.
\end{proof}








\begin{multicols}{2}[\section{Autres chapitres}]
\noindent
Préliminaires
\begin{enumerate}
\item \hyperref[introduction-section-phantom]{Introduction}
\item \hyperref[conventions-section-phantom]{Conventions}
\item \hyperref[sets-section-phantom]{Théorie des ensembles}
\item \hyperref[categories-section-phantom]{Catégories}
\item \hyperref[topology-section-phantom]{Topologie}
\item \hyperref[sheaves-section-phantom]{Faisceaux sur les espaces}
\item \hyperref[sites-section-phantom]{Sites et faisceaux}
\item \hyperref[stacks-section-phantom]{Champs}
\item \hyperref[fields-section-phantom]{Corps}
\item \hyperref[algebra-section-phantom]{Algèbre commutative}
\item \hyperref[brauer-section-phantom]{Groupes de Brauer}
\item \hyperref[homology-section-phantom]{Algèbre homologique}
\item \hyperref[derived-section-phantom]{Catégories dérivées}
\item \hyperref[simplicial-section-phantom]{Méthodes simpliciales}
\item \hyperref[more-algebra-section-phantom]{Compléments d'algèbre}
\item \hyperref[smoothing-section-phantom]{Lissage des morphismes d'anneaux}
\item \hyperref[modules-section-phantom]{Faisceaux de modules}
\item \hyperref[sites-modules-section-phantom]{Modules sur les sites}
\item \hyperref[injectives-section-phantom]{Objets injectifs}
\item \hyperref[cohomology-section-phantom]{Cohomologie des faisceaux}
\item \hyperref[sites-cohomology-section-phantom]{Cohomologie sur les sites}
\item \hyperref[dga-section-phantom]{Algèbre différentielle graduée}
\item \hyperref[dpa-section-phantom]{Algèbre à puissances divisées}
\item \hyperref[sdga-section-phantom]{Faisceaux différentiels gradués}
\item \hyperref[hypercovering-section-phantom]{Hyperrecouvrements}
\end{enumerate}
Schémas
\begin{enumerate}
\setcounter{enumi}{25}
\item \hyperref[schemes-section-phantom]{Schémas}
\item \hyperref[constructions-section-phantom]{Constructions de schémas}
\item \hyperref[properties-section-phantom]{Propriétés des schémas}
\item \hyperref[morphisms-section-phantom]{Morphismes de schémas}
\item \hyperref[coherent-section-phantom]{Cohomologie des schémas}
\item \hyperref[divisors-section-phantom]{Diviseurs}
\item \hyperref[limits-section-phantom]{Limites de schémas}
\item \hyperref[varieties-section-phantom]{Variétés}
\item \hyperref[topologies-section-phantom]{Topologies sur les schémas}
\item \hyperref[descent-section-phantom]{Descente}
\item \hyperref[perfect-section-phantom]{Catégories dérivées de schémas}
\item \hyperref[more-morphisms-section-phantom]{Compléments sur les morphismes}
\item \hyperref[flat-section-phantom]{Compléments sur la platitude}
\item \hyperref[groupoids-section-phantom]{Schémas en groupoïdes}
\item \hyperref[more-groupoids-section-phantom]{Compléments sur les schémas en groupoïdes}
\item \hyperref[etale-section-phantom]{Morphismes étales de schémas}
\end{enumerate}
Thèmes de théorie des schémas
\begin{enumerate}
\setcounter{enumi}{41}
\item \hyperref[chow-section-phantom]{Homologie de Chow}
\item \hyperref[intersection-section-phantom]{Théorie de l'intersection}
\item \hyperref[pic-section-phantom]{Schémas de Picard des courbes}
\item \hyperref[weil-section-phantom]{Théories cohomologiques de Weil}
\item \hyperref[adequate-section-phantom]{Modules adéquats}
\item \hyperref[dualizing-section-phantom]{Complexes dualisants}
\item \hyperref[duality-section-phantom]{Dualité pour les schémas}
\item \hyperref[discriminant-section-phantom]{Discriminants et différentes}
\item \hyperref[derham-section-phantom]{Cohomologie de de Rham}
\item \hyperref[local-cohomology-section-phantom]{Cohomologie locale}
\item \hyperref[algebraization-section-phantom]{Géométrie algébrique et formelle}
\item \hyperref[curves-section-phantom]{Courbes algébriques}
\item \hyperref[resolve-section-phantom]{Résolution des surfaces}
\item \hyperref[models-section-phantom]{Réduction semi-stable}
\item \hyperref[functors-section-phantom]{Foncteurs et morphismes}
\item \hyperref[equiv-section-phantom]{Catégories dérivées de variétés}
\item \hyperref[pione-section-phantom]{Groupes fondamentaux de schémas}
\item \hyperref[etale-cohomology-section-phantom]{Cohomologie étale}
\item \hyperref[crystalline-section-phantom]{Cohomologie cristalline}
\item \hyperref[proetale-section-phantom]{Cohomologie pro-étale}
\item \hyperref[relative-cycles-section-phantom]{Cycles relatifs}
\item \hyperref[more-etale-section-phantom]{Compléments de cohomologie étale}
\item \hyperref[trace-section-phantom]{La formule des traces}
\end{enumerate}
Espaces algébriques
\begin{enumerate}
\setcounter{enumi}{64}
\item \hyperref[spaces-section-phantom]{Espaces algébriques}
\item \hyperref[spaces-properties-section-phantom]{Propriétés des espaces algébriques}
\item \hyperref[spaces-morphisms-section-phantom]{Morphismes d'espaces algébriques}
\item \hyperref[decent-spaces-section-phantom]{Espaces algébriques décents}
\item \hyperref[spaces-cohomology-section-phantom]{Cohomologie des espaces algébriques}
\item \hyperref[spaces-limits-section-phantom]{Limites d'espaces algébriques}
\item \hyperref[spaces-divisors-section-phantom]{Diviseurs sur les espaces algébriques}
\item \hyperref[spaces-over-fields-section-phantom]{Espaces algébriques sur des corps}
\item \hyperref[spaces-topologies-section-phantom]{Topologies sur les espaces algébriques}
\item \hyperref[spaces-descent-section-phantom]{Descente et espaces algébriques}
\item \hyperref[spaces-perfect-section-phantom]{Catégories dérivées d'espaces}
\item \hyperref[spaces-more-morphisms-section-phantom]{Compléments sur les morphismes d'espaces}
\item \hyperref[spaces-flat-section-phantom]{Platitude sur les espaces algébriques}
\item \hyperref[spaces-groupoids-section-phantom]{Groupoïdes dans les espaces algébriques}
\item \hyperref[spaces-more-groupoids-section-phantom]{Compléments sur les groupoïdes dans les espaces}
\item \hyperref[bootstrap-section-phantom]{Amorçage}
\item \hyperref[spaces-pushouts-section-phantom]{Sommes amalgamées d'espaces algébriques}
\end{enumerate}
Thèmes de géométrie
\begin{enumerate}
\setcounter{enumi}{81}
\item \hyperref[spaces-chow-section-phantom]{Groupes de Chow des espaces}
\item \hyperref[groupoids-quotients-section-phantom]{Quotients de groupoïdes}
\item \hyperref[spaces-more-cohomology-section-phantom]{Compléments sur la cohomologie des espaces}
\item \hyperref[spaces-simplicial-section-phantom]{Espaces simpliciaux}
\item \hyperref[spaces-duality-section-phantom]{Dualité pour les espaces}
\item \hyperref[formal-spaces-section-phantom]{Espaces algébriques formels}
\item \hyperref[restricted-section-phantom]{Algébrisation des espaces formels}
\item \hyperref[spaces-resolve-section-phantom]{Résolution des surfaces revisitée}
\end{enumerate}
Théorie des déformations
\begin{enumerate}
\setcounter{enumi}{89}
\item \hyperref[formal-defos-section-phantom]{Théorie formelle des déformations}
\item \hyperref[defos-section-phantom]{Théorie des déformations}
\item \hyperref[cotangent-section-phantom]{Le complexe cotangent}
\item \hyperref[examples-defos-section-phantom]{Problèmes de déformation}
\end{enumerate}
Champs algébriques
\begin{enumerate}
\setcounter{enumi}{93}
\item \hyperref[algebraic-section-phantom]{Champs algébriques}
\item \hyperref[examples-stacks-section-phantom]{Exemples de champs}
\item \hyperref[stacks-sheaves-section-phantom]{Faisceaux sur les champs algébriques}
\item \hyperref[criteria-section-phantom]{Critères de représentabilité}
\item \hyperref[artin-section-phantom]{Axiomes d'Artin}
\item \hyperref[quot-section-phantom]{Espaces de Quot et de Hilbert}
\item \hyperref[stacks-properties-section-phantom]{Propriétés des champs algébriques}
\item \hyperref[stacks-morphisms-section-phantom]{Morphismes de champs algébriques}
\item \hyperref[stacks-limits-section-phantom]{Limites de champs algébriques}
\item \hyperref[stacks-cohomology-section-phantom]{Cohomologie des champs algébriques}
\item \hyperref[stacks-perfect-section-phantom]{Catégories dérivées de champs}
\item \hyperref[stacks-introduction-section-phantom]{Introduction aux champs algébriques}
\item \hyperref[stacks-more-morphisms-section-phantom]{Compléments sur les morphismes de champs}
\item \hyperref[stacks-geometry-section-phantom]{La géométrie des champs}
\end{enumerate}
Thèmes de théorie des espaces de modules
\begin{enumerate}
\setcounter{enumi}{107}
\item \hyperref[moduli-section-phantom]{Champs de modules}
\item \hyperref[moduli-curves-section-phantom]{Espaces de modules de courbes}
\end{enumerate}
Divers
\begin{enumerate}
\setcounter{enumi}{109}
\item \hyperref[examples-section-phantom]{Exemples}
\item \hyperref[exercises-section-phantom]{Exercices}
\item \hyperref[guide-section-phantom]{Guide bibliographique}
\item \hyperref[desirables-section-phantom]{Desiderata}
\item \hyperref[coding-section-phantom]{Style de programmation}
\item \hyperref[obsolete-section-phantom]{Obsolète}
\item \hyperref[fdl-section-phantom]{Licence de documentation libre GNU}
\item \hyperref[index-section-phantom]{Index généré automatiquement}
\end{enumerate}
\end{multicols}


\bibliography{my}
\bibliographystyle{amsalpha}

\end{document}
